% =====================================================================
%  Why Graph-Structured Systems Engaged by Finite Agents
%  Must Admit Operational Structure: A Conditional Derivation
%
%  Author: Kundai Farai Sachikonye
%          Technical University of Munich
% =====================================================================

\documentclass[11pt,twocolumn,a4paper]{article}

\usepackage[utf8]{inputenc}
\usepackage[T1]{fontenc}
\usepackage{lmodern}
\usepackage[a4paper,margin=2.4cm]{geometry}
\usepackage{amsmath,amssymb,amsthm,mathtools}
\usepackage{bm}
\usepackage{mathrsfs}
\usepackage{enumitem}
\usepackage{booktabs}
\usepackage{microtype}
\usepackage[round,sort&compress]{natbib}
\usepackage[colorlinks=true,linkcolor=blue!45!black,
            citecolor=blue!45!black,urlcolor=blue!45!black]{hyperref}
\usepackage{cleveref}

% ---- Theorem environments ----
\newtheoremstyle{thmstyle}{\topsep}{\topsep}{\itshape}{}{\bfseries}{.}{.5em}{}
\theoremstyle{thmstyle}
\newtheorem{theorem}{Theorem}[section]
\newtheorem{lemma}[theorem]{Lemma}
\newtheorem{proposition}[theorem]{Proposition}
\newtheorem{corollary}[theorem]{Corollary}
\theoremstyle{definition}
\newtheorem{definition}[theorem]{Definition}
\newtheorem{axiom}{Premise}
\newtheorem{principle}[theorem]{Principle}
\theoremstyle{remark}
\newtheorem{remark}[theorem]{Remark}
\newtheorem{example}[theorem]{Example}

\crefname{axiom}{Premise}{Premises}
\Crefname{axiom}{Premise}{Premises}

% ---- Notation ----
\newcommand{\R}{\mathbb{R}}
\newcommand{\N}{\mathbb{N}}
\newcommand{\Z}{\mathbb{Z}}
\newcommand{\Zp}{\mathbb{Z}_{\geq 0}}
\newcommand{\Graph}{\mathcal{G}}
\newcommand{\Vrt}{\mathsf{V}}
\newcommand{\Edg}{\mathsf{E}}
\newcommand{\Agent}{\mathcal{A}}
\newcommand{\Op}{\mathsf{Op}}
\newcommand{\Res}{\rho}
\newcommand{\floor}{\beta}
\newcommand{\OpStr}{\mathfrak{O}}

\title{\textbf{On the Strict Necessity of Finite Weighted Graph Structures In Bounded Resolvable Space:\\ Conditional Derivation of Operational Structure Requirements for Finite Agents in Bounded Systems}}

\author{%
Kundai Farai Sachikonye\\
\small Technical University of Munich\\
\small AIMe Registry for Artificial Intelligence\\
\small \texttt{kundai.sachikonye@wzw.tum.de}}

\date{\today}

% =====================================================================
\begin{document}
\maketitle

\begin{abstract}
\noindent
We give a conditional derivation: granted two premises---that the
engaging agents are finite, and that the engaged system is a finite
graph---we show that the system necessarily admits an
\emph{operational structure} in a precise sense. The argument is
constructive. We define a finite graph-structured system and a finite
agent, derive that any distinguishability relation an agent can
realise has a strictly positive floor on edge measure, prove that
this floor forces a residue on every agent--graph interaction, and
show that the realised acts form a small category whose hom-sets
carry a preorder by residue. This category, presented with its
residue preorder and a budget-bounded notion of realisability, is
what we call the system's operational structure. The structure is
not chosen; it is the unique closure of the agent's realisable
interactions under the constraints imposed by finiteness on both
sides. We classify operations into four canonical types---algebraic,
dynamical, computational, and internal---and show every realisable
operation is a combination of these. The derivation makes no
metaphysical commitments and applies wherever the two premises
hold; the framework is silent outside that scope. We close by
stating the narrow falsifiability of the main theorem: it can be
refuted only by exhibiting a finite agent realising
distinguishability with zero floor on a finite graph, which we prove
is structurally impossible. The paper is self-contained; no result
of the author is cited; the literature references are to standard
mathematical lemmas.
\end{abstract}

\medskip
\noindent\textbf{Keywords:} finite agents; finite graphs; operational
structure; distinguishability floor; residue propagation; small
category; partial composition; residue preorder; conditional
derivation.

\tableofcontents

% =====================================================================
\section{Introduction}
\label{sec:intro}
% =====================================================================

\subsection{The conditional question}

The question this paper addresses is short:
\begin{quote}\itshape
On the condition that the engaging agents are finite and the engaged
system is a finite graph, what structure of operations must the
system admit?
\end{quote}
We answer: a specific operational structure---a small category
with a residue preorder, a positive floor, and a budget-bounded
realisability marker---derivable from the two premises alone, with
no further hypotheses.

The result is conditional throughout. We do not assert that all
agents are finite, or that all systems are graphs, or that systems
are fundamentally relational. We assert only that \emph{if} both
conditions hold, the operational structure of
\cref{thm:opstructure} is forced. The framework is silent outside
that scope: where either premise fails, the derivation does not
apply and makes no prediction.

\subsection{Why this matters}

Engineering practice routinely treats operational structures---
function composition, command pipelines, state-machine transitions,
query algebras---as design choices to be selected for convenience.
The argument here shows that for the dominant case (finite agents
engaging finite graphs, which covers essentially all of practical
software, embedded systems, sensor networks, knowledge graphs, and
finite databases) the operational structure is \emph{not} a design
choice but a derivable necessity. The class of admissible structures
is forced; only the labelling within the class is free.

This converts a large family of engineering decisions from
\emph{which structure to invent} into \emph{which structure has the
two premises already determined}. The latter is a discovery task
with a definite answer; the former is a design task with no
objective ground.

\subsection{Method}

The derivation is constructive. We fix the two premises in
\cref{sec:premises}, derive three consequences in
\cref{sec:consequences} (distinguishability floor, residue
propagation, small-category structure under partial composition),
prove the main theorem in \cref{sec:main}, classify the resulting
operations in \cref{sec:classification}, and locate the
falsifiability of the argument in \cref{sec:falsifiability}. Every
step is a finite combinatorial, measure-theoretic, or
category-theoretic claim with a self-contained proof; no result of
the author is cited.

\subsection{Relation to existing frameworks}

The argument touches finite graph theory \citep{diestel2017,
bollobas1998}, finite-state automata \citep{hopcroft2007}, measure
theory of finite partitions \citep{halmos1950}, category theory and
partial composition \citep{maclane1998}, semigroup and monoid theory
\citep{howie1995}, and order theory on finite lattices
\citep{daveypriestley2002}. We use only the textbook content of
each field. The novel contribution is the closure argument
(\cref{thm:opstructure}) that forces a small-category operational
structure with floor, residue, and preorder from the conjunction of
the two premises; the constituent pieces are standard, but their
conjunction has not been packaged as a forced structural result, to
the author's knowledge.

% =====================================================================
\section{Three Premises}
\label{sec:premises}
% =====================================================================

We fix three premises. The first two are operational conditions on
the agent and the system; the third is an order condition on the
ambient whole that the agent's graph models. None is a metaphysical
hypothesis. The third premise carries the load that the original
draft tried to put on the floor by stipulation: it forces the floor
as a derived consequence rather than a posit, and it converts the
partiality of composition, the cell-structure of the residue
preorder, and the four-type taxonomy into theorems rather than
defects.

\subsection{Finite graph-structured system}

\begin{definition}[Finite graph-structured system]
\label{def:graph}
A \emph{finite graph-structured system} is a triple $\Graph =
(\Vrt, \Edg, w)$ where $\Vrt$ is a finite non-empty set of
\emph{vertices}, $\Edg \subseteq \Vrt \times \Vrt$ is a finite set
of \emph{edges}, and $w: \Edg \to \R_{\geq 0}$ is an \emph{edge
measure} assigning to each edge a non-negative real weight. We do
not require edges to be undirected or distinct; multigraphs are
admitted.
\end{definition}

\begin{axiom}[Finite graph hypothesis]
\label{ax:graph}
The engaged system $\Graph = (\Vrt, \Edg, w)$ is a finite
graph-structured system: $|\Vrt| < \infty$, $|\Edg| < \infty$, and
$w$ takes finite values throughout.
\end{axiom}

The graph need not be connected, simple, or planar. It need not be
metric, in the sense that $w$ need not satisfy the triangle
inequality. We require only that it has a finite number of
vertices, a finite number of edges, and a finite-valued non-negative
measure on its edges.

\begin{remark}[Why graphs]
\label{rem:why_graphs}
We pick graphs as the primitive structure because they are the
minimal combinatorial encoding of relational data: a set of entities
together with the relations between them. Any system describable as
``entities standing in relations'' is encodable as a graph, possibly
a multigraph with labelled edges. Trees, finite-state machines,
knowledge bases, relational databases, dependency graphs, and finite
hypergraphs all reduce to or extend this minimal form. The choice
restricts neither the scope of application nor the generality of the
result.
\end{remark}

\subsection{Finite agent}

\begin{definition}[Distinguishability relation]
\label{def:distrel}
A \emph{distinguishability relation} on $\Graph$ is a binary
relation $\sim \subseteq \Vrt \times \Vrt$ such that for each pair
$(u, v)$, the agent has either established $u \sim v$ (the pair is
treated as a single entity), or $u \not\sim v$ (the pair is treated
as distinct entities), but not both.
\end{definition}

\begin{definition}[Finite agent]
\label{def:agent}
A \emph{finite agent} $\Agent$ on $\Graph$ is specified by:
\begin{enumerate}[label=(F\arabic*),leftmargin=2.6em,nosep]
\item\label{F:fin_state} A finite internal state space $S_\Agent$
with $|S_\Agent| < \infty$, encoding all distinctions the agent can
hold at once.
\item\label{F:fin_descr} A finite descriptive scheme: each
distinguishability decision is recorded by an injection from the
$\sim$-classes into $S_\Agent$, so the number of distinct
equivalence classes the agent can simultaneously hold is at most
$|S_\Agent|$.
\item\label{F:fin_cost} A finite cost budget: each act of
distinguishing a pair $(u,v)$ as $u \not\sim v$ has a cost
$c(u,v) > 0$ bounded below by a positive constant $c_0$, and the
agent's total expenditure over its lifetime is bounded by a finite
budget $C_{\Agent} < \infty$.
\end{enumerate}
\end{definition}

\begin{axiom}[Finite agent hypothesis]
\label{ax:agent}
Any agent engaging $\Graph$ is a finite agent in the sense of
\cref{def:agent}: it has finite internal state, finite descriptive
scheme, and finite cost budget.
\end{axiom}

We make no further assumptions about the agent's architecture: it
may be a human, a machine, a sensor network, a software process, a
finite collective. The three finiteness conditions of
\cref{def:agent} are the substantive content of \cref{ax:agent}.

\begin{figure*}[!htbp]
    \centering
    \includegraphics[width=\textwidth]{figures/panel_02_graph_floor.png}
    \caption{\textbf{Graph-Level Floor on Realised Distinguishability Paths.}
    (\textbf{A})~Maximum edge weight on each realised path on the linear $y$-axis versus graph-level floor $\beta$ on the linear $x$-axis, $1{,}992$ realised distinguishability acts across $40$ random graphs ($|V|=15$, edge probability $0.3$). Steelblue points are realised acts; the neutral dashed line marks $\max\text{-edge} = \beta$. All $1{,}992$ points lie on or above the line; zero violations. Empirical content of Theorem~\ref{thm:floor} (Floor on Edge Measure): every realised distinguishability path contains an edge of measure $\geq \beta$.
    (\textbf{B})~Residue $\rho$ on the linear $y$-axis versus graph-level floor $\beta$ on the linear $x$-axis for the same $1{,}992$ acts. The neutral dashed line marks $\rho = \beta$. All points lie above the line, with residues spreading upward as longer paths are traversed. This is the immediate consequence of (A): subadditivity of residue along realised paths gives $\rho \geq \max\text{-edge} \geq \beta$.
    (\textbf{C})~Residue-floor ratio $\rho/\beta$ on the log $x$-axis ($10^{-1}$ to $10^{2}$) versus count on the linear $y$-axis. The crimson dashed reference at $\rho/\beta = 1$ marks the floor exactly; every observed value lies at or above $1$, with the bulk of the distribution at $\rho/\beta \in [1, 10]$. The histogram visualises the headroom realised acts carry above the floor.
    (\textbf{D})~Three-dimensional cloud of $(\beta,\,\rho,\,|V|)$, viridis colourmap by $\rho/\beta$, viewing angle elevation $30^{\circ}$ azimuth $-60^{\circ}$. Each point is one realised act; the cloud lies entirely above the diagonal plane $\rho = \beta$ at every $|V|$ tested. Geometric realisation of Theorem~\ref{thm:floor}: the floor is a graph property and applies uniformly across vertex counts.}
    \label{fig:panel-graph-floor}
\end{figure*}

\subsection{Non-completable whole}

The agent engages a finite graph, but the graph is a
\emph{model}---a finite individuated structure within an ambient
whole that the agent's model does not exhaust. We name the
ambient whole and require that it is not exhaustible by finite
stages of modelling.

\begin{definition}[Whole; medium]
\label{def:medium}
The \emph{whole} engaged by an agent is the totality
$\mathcal{M}$ within which individuation occurs. The graph
$\Graph = (\Vrt, \Edg, w)$ is a finitely-individuated model of
$\mathcal{M}$. The \emph{medium} is $\mathcal{M}$ in its role as
the reference against which the vertices of $\Graph$ are
individuated: the totality simultaneously in contact with every
vertex, such that the cut weight of a subset $U \subseteq \Vrt$ is
the cost of separating $U$ from $\mathcal{M}\setminus U$.
\end{definition}

\begin{definition}[Non-completable whole]
\label{def:noncompletable}
The whole $\mathcal{M}$ is \emph{non-completable} if no finite
stage of individuation exhausts it: for every finite individuation
$\Graph$, there is at least one further distinction the agent could
record but has not, i.e. a refinement of $\Graph$ that the agent's
descriptive scheme can in principle extend to. This is an order
condition (no terminal stage), not a cardinality posit on
$\mathcal{M}$.
\end{definition}

\begin{axiom}[Non-completability hypothesis]
\label{ax:noncompletable}
The whole $\mathcal{M}$ the agent engages is non-completable in the
sense of \cref{def:noncompletable}.
\end{axiom}

\Cref{ax:noncompletable} is not the claim that the world is
infinite. It is the operational claim that the agent's modelling
process does not terminate at a finite stage: there is always at
least one more distinction available to draw beyond the current
graph. For any realised engineering system this is observably the
case---a sensor of finite resolution always admits a finer one; a
classifier of finite vocabulary always admits a wider one; a query
language of finite expressiveness always admits an extension. The
operational content of \cref{ax:noncompletable} is exactly the
practical fact that finite-agent modelling is open-ended in the
direction of refinement.

\begin{remark}[The third premise is what carries the floor]
\label{rem:third_carries}
In the original draft the positive floor $\floor$ on edge weights
was stipulated as part of \cref{ax:graph}. Under
\cref{ax:noncompletable}, $\floor > 0$ is no longer a stipulation
but a theorem (\cref{thm:floor_from_infinitude} below): the floor
is the trace of the non-completable whole on its finitely-modelled
parts. With the floor derived, the remaining structural claims
become theorems about the agent's engagement with the
non-completable whole rather than artefacts of a stipulated bound.
\end{remark}

\subsection{What the premises do not say}

\Cref{ax:graph,ax:agent} are sharp where they need to be sharp and
silent elsewhere. We list three things the premises explicitly do
not say.

\begin{enumerate}[nosep,leftmargin=*]
\item They do not say the world is finite. Reality at large is
permitted to be infinite; the premises restrict only the system the
agent engages and the agent itself, both of which are finite.
\item They do not say the graph encodes everything about the system.
A graph is a \emph{model} of the system the agent uses; what the
system is in itself, outside the agent's modelling, is not
addressed.
\item They do not say agents must be physical. The argument applies
to any agent satisfying~\ref{F:fin_state}--\ref{F:fin_cost},
whatever its substrate.
\end{enumerate}

\subsection{Scope}

The conjunction of the two premises covers the dominant case of
engineering practice: software systems, embedded controllers,
sensor networks, knowledge bases, finite databases, finite-state
protocols, and finite communication graphs. It does not cover
putative infinite agents (perfect oracles, unbounded computation) or
putative infinite systems (continuous fields engaged directly,
without finite discretisation). Where finite discretisation is
imposed---which is the case in all realised engineering
systems---the premises apply.

% =====================================================================
\section{Consequences of the Three Premises}
\label{sec:consequences}
% =====================================================================

We derive four consequences of the conjunction of \cref{ax:graph},
\cref{ax:agent}, and \cref{ax:noncompletable}. The first is the
Floor from Infinitude theorem: the positive floor on edge weights
that the original draft posited is now derived from
non-completability. The remaining three---distinguishability floor,
residue propagation, and small-category structure---are then
recovered under the derived floor and the agent's finiteness. Each
consequence is forced by the premises alone; no additional
hypothesis is invoked.

\subsection{The Floor from Infinitude}

\begin{theorem}[Floor from Infinitude]
\label{thm:floor_from_infinitude}
Let $\Graph$ be a finitely-individuated model of a non-completable
whole $\mathcal{M}$ (\cref{ax:noncompletable}). For every realisable
distinguishability act $u \not\sim v$ by a finite agent
(\cref{ax:agent}), the agent's individuation of $u$ from
$\mathcal{M}\setminus\{u\}$ via the path through $\Graph$ deposits a
strictly positive residual: no realisable individuation has zero
residue. Equivalently, the cost of any realised separator in
$\Graph$ is bounded below by a strictly positive constant.
\end{theorem}

\begin{proof}
A realisable distinguishability act $u \not\sim v$ individuates $u$
from the rest of the agent's model by depositing measure on a
separator $\Sigma(u, v) \subseteq \Edg$ (the edges of the path the
agent uses). To individuate $u$ \emph{completely} would be to
finish the comparison of $u$ against the entirety of
$\mathcal{M}\setminus\{u\}$. By \cref{ax:noncompletable},
$\mathcal{M}\setminus\{u\}$ is non-completable: no finite stage of
individuation exhausts it. Hence the comparison cannot finish at
any finite stage. Therefore the separator $\Sigma(u, v)$ cannot
have zero measure: were $w(\Sigma) = 0$, the agent's descriptive
scheme would record no measure differential between $u$ and the
rest, and $u$ would not be individuated. Since the act is realised,
the separator carries strictly positive measure, and the infimum of
such positive measures over the finite edge set $\Edg$ is itself a
strictly positive constant
$\floor := \min\{w(e) : e \in \Edg,\, w(e) > 0\} > 0$. $\qed$
\end{proof}

\begin{corollary}[The floor is the trace of the whole on its
parts]
\label{cor:floor_trace}
$\floor$ is not an independent posit but the structural shadow that
$\mathcal{M}$'s non-completability casts on each of its finite
modelled parts. Were $\mathcal{M}$ completable (a finite whole),
the comparison of any part against the rest could finish, the
separator could shrink to zero measure, and $\floor = 0$ would be
admissible. Because $\mathcal{M}$ is non-completable
(\cref{ax:noncompletable}), the floor is positive.
\end{corollary}

\begin{proof}
Immediate from \cref{thm:floor_from_infinitude}: $\floor > 0$ is
forced by non-completability and would vanish in the completable
case. The constant $\floor$ is the irreducible identification
residual of the agent's finite individuating acts against the
non-completable whole. $\qed$
\end{proof}

\begin{figure*}[!htbp]
    \centering
    \includegraphics[width=\textwidth]{figures/panel_01_floor_from_infinitude.png}
    \caption{\textbf{Floor from Infinitude: the floor $\beta$ is the trace of the non-completable whole on its finitely-modelled parts.}
    (\textbf{A})~Floor $\beta$ on a log $y$-axis ($10^{-2}$ to $10^{0}$) versus refinement stage $n$ on the linear $x$-axis ($0$ to $11$) across twelve stages of a non-completable whole $\mathcal{M}$. Steelblue circles connected by a line trace $\beta(n)$; the floor decreases monotonically from $\beta=0.5$ at stage $0$ to $\beta\approx 0.077$ at stage $11$ yet remains strictly positive at every tested stage. This is the empirical content of Theorem~\ref{thm:floor_from_infinitude} (Floor from Infinitude): finite individuation against a non-completable whole forces a strictly positive residual at every finite stage.
    (\textbf{B})~Decay shape: measured $\beta(n)$ (steelblue circles) versus the theoretical refinement law $\beta_{\text{theory}}(n)=1/(n+2)$ (crimson solid reference) on the linear $y$-axis ($0$ to $0.5$) versus stage $n$ on the linear $x$-axis ($0$ to $11$). The measured values follow the prediction to floating-point precision, confirming that the refinement law in the harness reproduces the trace structure of Corollary~\ref{cor:floor_trace}.
    (\textbf{C})~Graph size growth: $|V|$ (steelblue circles) and $|E|$ (teal squares) on the linear $y$-axis (up to $400$) versus stage $n$ on the linear $x$-axis ($0$ to $11$). Vertex count grows linearly ($|V|=5+3n$); edge count grows quadratically. The non-completable whole is unbounded in extension at every finite stage, capturing Definition~\ref{def:noncompletable} (no terminal stage).
    (\textbf{D})~Three-dimensional scatter of edge-weight strata: $\log_{10} w$ on the vertical axis as a function of stage $n$ on the linear $x$-axis ($0$ to $7$) and normalised edge rank on the linear $y$-axis ($0$ to $1$), viridis colourmap by stage, viewing angle elevation $30^{\circ}$ azimuth $-60^{\circ}$. Each stage contributes a strictly finer family of edges, visible as descending viridis-coloured fans; the floor at each stage is the bottom of its fan and stays bounded below. Geometric realisation of Theorem~\ref{thm:floor_from_infinitude}: the floor is the bottom of an indefinitely extensible spectrum, never reaching zero.}
    \label{fig:panel-floor-from-infinitude}
\end{figure*}

\subsection{The distinguishability floor}

\begin{theorem}[Floor on edge measure for realisable
distinguishability]
\label{thm:floor}
Let $\Graph$ satisfy \cref{ax:graph} and $\Agent$ satisfy
\cref{ax:agent}. Define the graph-level floor
\[
  \floor := \min\{w(e) : e \in \Edg,\; w(e) > 0\}.
\]
This is a property of $\Graph$ alone, well-defined whenever
$\Graph$ contains at least one positive-measure edge. Then
$\floor > 0$, and for every pair $(u,v) \in \Vrt \times \Vrt$ that
$\Agent$ can realise as $u \not\sim v$, the agent's
distinguishability path through $\Graph$ contains an edge of measure
$\geq \floor$.
\end{theorem}

\begin{proof}
\emph{Positivity of $\floor$.} The set
$W^{+} := \{w(e) : e \in \Edg,\; w(e) > 0\}$ is a finite (by
$|\Edg| < \infty$ in \cref{ax:graph}) subset of $\R_{> 0}$. If
$W^{+}$ is non-empty, $\floor = \min W^{+} > 0$. If $W^{+}$ is
empty, every edge has measure $0$, no path can have positive
total measure, and \cref{F:fin_descr} prevents the agent from
recording any non-trivial distinguishability (a null path gives the
agent no measure differential to record); in this degenerate case
no realisable distinguishability acts exist and the claim holds
vacuously.

\emph{Realised paths contain a $\geq \floor$ edge.} Assume the
non-degenerate case, so $\floor > 0$. Let $(u,v)$ be realised by
$\Agent$ as $u \not\sim v$. By \cref{F:fin_descr}, the agent
records the distinction by some path
$P_\Agent(u,v) = (e_1, e_2, \ldots, e_k)$ in $\Graph$. If every
edge of $P_\Agent(u,v)$ had measure $0$, the path would have total
measure $0$, the agent's descriptive scheme would have no measure
differential to encode the distinction, and the distinguishability
would not be realised---contradicting the assumption. Hence at
least one edge $e_i$ of $P_\Agent(u,v)$ has $w(e_i) > 0$, and
therefore $w(e_i) \geq \floor$ by the definition of $\floor$ as
the minimum positive edge measure. $\qed$
\end{proof}

\begin{remark}[The floor is a graph property, not an agent
property]
\label{rem:floor_graph_property}
$\floor$ depends only on $\Graph$, not on $\Agent$. Different
finite agents engaging the same graph share the same floor.
$\Agent$ enters only in the second clause, where realisability
guarantees that at least one edge above the floor lies on every
realised distinguishability path.
\end{remark}

\begin{remark}[Floor is a finiteness floor]
\label{rem:floor_is_finiteness}
The floor $\floor$ is not a metaphysical lower bound on what edges
``can be.'' It is the lower bound forced by the finiteness of
$\Edg$ and the requirement that distinguishability be realisable.
With a finite edge set having positive-measure edges, the minimum
positive measure is itself positive. This is a combinatorial
observation, not a substantive claim about the structure of
reality.
\end{remark}

\subsection{Residue: every distinguishability act deposits content}

\begin{definition}[Residue]
\label{def:residue}
The \emph{residue} $\Res(u,v)$ of a distinguishability act
$u \not\sim v$ realised by $\Agent$ is the sum of edge measures
along the path the agent uses:
\[
  \Res(u,v) = \sum_{e \in \mathrm{path}_\Agent(u,v)} w(e).
\]
By \cref{thm:floor}, $\Res(u,v) \geq \floor > 0$ for every
realisable distinguishability act.
\end{definition}

\begin{theorem}[Residues are additive under composition]
\label{thm:residue_additive}
Let $u \not\sim v$ and $v \not\sim w$ be two distinguishability
acts realised by $\Agent$. The composed act $u \not\sim w$ (if
realisable) has residue:
\[
  \Res(u, w) \leq \Res(u, v) + \Res(v, w).
\]
Equality holds when the path the agent uses for the composed act is
the concatenation of the individual paths.
\end{theorem}

\begin{proof}
The residue of the composed act is the sum of edge measures along
the agent's path from $u$ to $w$. By the triangle inequality for
sums of non-negative numbers, the shortest such path has total
measure no greater than the sum of any concatenated route. If the
agent uses the concatenation $u \to v \to w$, equality holds;
otherwise the agent has found a shorter path and the inequality is
strict. $\qed$
\end{proof}



\begin{figure*}[!htbp]
    \centering
    \includegraphics[width=\textwidth]{figures/panel_03_residue_propagation.png}
    \caption{\textbf{Residue Subadditivity and Cumulative Monotone Growth.}
    (\textbf{A})~Composite residue $\rho(u,w)$ on the linear $y$-axis versus the sum $\rho(u,v)+\rho(v,w)$ on the linear $x$-axis for $1{,}174$ random triples across $30$ random graphs ($|V|=15$, edge probability $0.4$). Steelblue points lie at or below the neutral diagonal $y=x$, with $1{,}043$ trials strictly below (the agent found a shorter route than the $u\to v\to w$ concatenation) and $131$ trials exactly on the line (concatenation was the shortest). Zero subadditivity violations. Empirical content of Theorem~\ref{thm:residue_additive}: $\rho(u,w) \leq \rho(u,v) + \rho(v,w)$.
    (\textbf{B})~Cumulative residue along $8$ independent realised walks on the linear $y$-axis ($0$ to $\sim 300$) versus act count $n$ on the linear $x$-axis ($0$ to $80$). Each steelblue curve is a separate walk; the corresponding crimson floor line $n\cdot\beta$ lies below it at every point. All eight cumulative trajectories are monotone non-decreasing, with $\sum\rho_i \geq n\beta$ holding throughout. Empirical content of Corollary~\ref{cor:res_accumulate}.
    (\textbf{C})~Subadditivity gap $\rho(u,v)+\rho(v,w)-\rho(u,w)$ on the linear $x$-axis (up to $\sim 18$) versus count on the linear $y$-axis. The neutral vertical reference at $0$ marks the equality case; the distribution lies entirely in the non-negative half-line. The mode is at small positive values, with a heavy tail of triples where the agent found a substantially shorter composite path.
    (\textbf{D})~Three-dimensional composition-triple cloud: $\rho(u,v)$ on the linear $x$-axis, $\rho(v,w)$ on the linear $y$-axis, $\rho(u,w)$ on the vertical axis, viridis colourmap by subadditivity gap, viewing angle elevation $30^{\circ}$ azimuth $-60^{\circ}$. The cloud lies below the diagonal plane $\rho(u,w) = \rho(u,v) + \rho(v,w)$; the subadditive structure is the geometric content of Theorem~\ref{thm:residue_additive}.}
    \label{fig:panel-residue-propagation}
\end{figure*}

\begin{corollary}[Residue accumulates and does not refund]
\label{cor:res_accumulate}
For any sequence of distinguishability acts
$u_0 \not\sim u_1 \not\sim u_2 \not\sim \cdots \not\sim u_n$, the
total residue deposited satisfies:
\[
  \sum_{i=0}^{n-1} \Res(u_i, u_{i+1}) \geq n\floor > 0.
\]
The cumulative residue strictly grows with the number of acts and
is bounded below by $n\floor$. No act reduces the cumulative total;
the sum is monotone non-decreasing in $n$.
\end{corollary}

\begin{proof}
By \cref{thm:floor}, each $\Res(u_i, u_{i+1}) \geq \floor$. Summing
$n$ non-negative terms each $\geq \floor$ gives $\geq n\floor$.
Monotonicity in $n$ is immediate from the non-negativity of each
term. $\qed$
\end{proof}

\subsection{Closure: distinguishability acts form a small category}

Composition of distinguishability acts is intrinsically
\emph{partial}: $(u,v)$ and $(v,w)$ both being realised by $\Agent$
does not imply that $(u,w)$ is realised. Realisability depends on
the agent's path-choices and on the budget $C_\Agent$. The
finite-budget premise \cref{F:fin_cost} actively limits which
compositions are realised, so a total composition law cannot be
assumed.

The honest object is therefore a \emph{small category}, not a
monoid. In a small category, composition is partial in exactly
this controlled sense: a morphism $g: u \to v$ composes with a
morphism $f: v \to w$ to give $f \circ g: u \to w$, but morphisms
with mismatched source/target do not compose. We construct the
category and verify the axioms.

\begin{definition}[The distinguishability category]
\label{def:dcat}
The \emph{distinguishability category}
$\mathbf{C}_{\Agent,\Graph}$ has:
\begin{itemize}[nosep,leftmargin=*]
\item \emph{Objects}: the vertices $\Vrt$ of $\Graph$.
\item \emph{Morphisms}: a morphism $u \to v$ is either the
identity $\mathrm{id}_u$ (with residue $0$), or a realised
distinguishability act $(u,v) \in \mathcal{D}_\Agent$ together with
the agent's path $\mathrm{path}_\Agent(u,v)$ used to realise it.
\item \emph{Composition}: for morphisms $g: u \to v$ and
$h: v \to w$, the composite $h \circ g: u \to w$ is defined iff
$\Agent$ realises a distinguishability act $(u,w)$ whose path is
recordable from the agent's record of $g$ and $h$. When defined,
$h \circ g = (u, w, \mathrm{path}_\Agent(u,w))$.
\item \emph{Identities}: for each $u \in \Vrt$, the identity
$\mathrm{id}_u : u \to u$ is the trivial act with empty path and
residue $0$.
\end{itemize}
\end{definition}

\begin{lemma}[Identity laws]
\label{lem:identity}
For every realised morphism $g: u \to v$,
$g \circ \mathrm{id}_u = g$ and $\mathrm{id}_v \circ g = g$.
\end{lemma}

\begin{proof}
The composites $g \circ \mathrm{id}_u$ and $\mathrm{id}_v \circ g$
have source $u$, target $v$, residue $\Res(g) + 0 = \Res(g)$, and
path $\mathrm{path}_\Agent(u,v)$. Both are the realised morphism
$g$. The identity composites are realisable because they require
no agent effort beyond what $g$ itself already required. $\qed$
\end{proof}

\begin{lemma}[Associativity where defined]
\label{lem:assoc}
For realised morphisms $g_1: u \to v$, $g_2: v \to w$,
$g_3: w \to x$ such that both $(g_3 \circ g_2) \circ g_1$ and
$g_3 \circ (g_2 \circ g_1)$ are defined,
\[
  (g_3 \circ g_2) \circ g_1 \;=\; g_3 \circ (g_2 \circ g_1).
\]
\end{lemma}

\begin{proof}
Both composites have source $u$ and target $x$. By
\cref{F:fin_descr}, the agent records a morphism by its
endpoints and its realised path; if both composites are defined,
the agent has recorded a realised morphism $u \to x$, and the
record is unique up to the agent's path-choice. The agent's
path-choice for the composite is a function only of the endpoints
and the available realised sub-morphisms, which both bracketings
present identically. Hence the two composites are the same
morphism. $\qed$
\end{proof}

\begin{theorem}[Realised acts form a small category]
\label{thm:category}
Under \cref{ax:graph} and \cref{ax:agent},
$\mathbf{C}_{\Agent,\Graph}$ is a small category: objects form a
finite set ($|\Vrt| < \infty$), morphisms form a finite set
(bounded by $|\Vrt| + |\mathcal{D}_\Agent| \leq |\Vrt| +
|S_\Agent|^2 < \infty$), composition is partial and associative
where defined (\cref{lem:assoc}), and identities exist
(\cref{lem:identity}).
\end{theorem}

\begin{proof}
Finiteness of objects and morphisms is immediate from the premises.
\cref{lem:identity,lem:assoc} provide the identity and associativity
axioms in the partial-composition setting. $\qed$
\end{proof}

\begin{theorem}[Partiality of composition is forced by
non-completability]
\label{thm:partiality_forced}
Under \cref{ax:graph}, \cref{ax:agent}, and
\cref{ax:noncompletable}, composition in
$\mathbf{C}_{\Agent,\Graph}$ is genuinely partial: there exist
agent--graph configurations in which $(u, v)$ and $(v, w)$ are both
realised but $(u, w)$ is not. Partiality is not a contingent
artefact of budget exhaustion; it is the structural consequence of
modelling a non-completable whole with a finite agent.
\end{theorem}



\begin{proof}
By \cref{thm:floor_from_infinitude}, every realised
distinguishability act deposits residue $\geq \floor > 0$. The
agent's individuating capacity over its lifetime is the cumulative
residue it can deposit, bounded above by what the finite budget
$C_\Agent$ permits at unit cost $c_0$:
$\#\{\text{realised acts}\} \leq C_\Agent / c_0$.
Composition of $(u, v)$ and $(v, w)$ into $(u, w)$ is realised only
when the agent can record the composite act---which requires either
that the composite path already lies in the agent's record, or that
the agent has remaining capacity to traverse it.

Now apply \cref{ax:noncompletable}: the whole $\mathcal M$ is
non-completable, so the lattice of pairs $(u, w)$ the agent
\emph{could} individuate exceeds any finite stage of realisations.
For any finite agent there are pairs $(u, v), (v, w)$ both realised
whose composite $(u, w)$ has not been realised at this stage---the
non-completability of $\mathcal M$ guarantees that finite
individuation is always strictly extensible, so the agent's record
is always a proper sub-lattice of the closure under composition.

Hence composition is partial, structurally: the gap between the
realised category and its composition closure is a non-empty
sub-lattice for every finite agent in a non-completable whole, and
collapsing the gap would require the agent to have exhausted the
whole (forbidden by \cref{ax:noncompletable}). $\qed$
\end{proof}

\begin{remark}[Small category is the honest structural object]
\label{rem:why_category}
\Cref{thm:partiality_forced} promotes partial composition from a
contingent feature of budget exhaustion to a forced consequence of
non-completability. The honest structural object is a small
category \citep{maclane1998}, not a monoid. The category framework
is exactly what the conjunction of the three premises supports: a
finite collection of morphisms (the realised individuations), an
identity on each object, an associative partial composition
(\cref{lem:assoc}), with the gap from totality not a defect but the
structural shadow of the non-completable whole. A small category in
which every pair of composable morphisms does compose corresponds
exactly to a completable whole (the agent has exhausted what is
individuable); under \cref{ax:noncompletable} this case is
excluded.
\end{remark}

\begin{remark}[Free monoid envelope]
\label{rem:free_monoid}
The path category over $\mathbf{C}_{\Agent,\Graph}$---formal
sequences of composable morphisms, taken modulo the realised
composites---is a free monoid on the morphisms of
$\mathbf{C}_{\Agent,\Graph}$, modulo the relations imposed by
realisation. This is a derived construction, useful for some
applications, but not the primitive structure forced by the
premises. The primitive structure is the small category.
\end{remark}

% --- Hom-set sizes ---

\begin{proposition}[Hom-set bound]
\label{prop:hom_bound}
For any $u, v \in \Vrt$, the hom-set
$\mathrm{Hom}_{\mathbf{C}_{\Agent,\Graph}}(u, v)$ contains at most
one realised non-identity morphism plus (if $u = v$) the identity,
totalling at most two elements.
\end{proposition}

\begin{proof}
A realised distinguishability act $(u, v)$ is the assertion
$u \not\sim v$; the agent records this assertion at most once in
its descriptive scheme (it is either established or not). If the
agent's path-choice changes during a run, the agent's record is
updated, but the morphism remains the unique realised act between
$u$ and $v$. The identity on $u$ exists when $u = v$. $\qed$
\end{proof}

\begin{figure*}[!htbp]
    \centering
    \includegraphics[width=\textwidth]{figures/panel_04_category.png}
    \caption{\textbf{Small-Category Structure: Hom-Set Bound, Identity Laws, and Morphism Counts.}
    (\textbf{A})~Hom-set size distribution on a log $y$-axis ($10^{0}$ to $10^{4}$) versus hom-set size on the linear $x$-axis ($\{0, 1, 2\}$). Across $9{,}000$ inspected hom-sets in $50$ random graphs, no hom-set has size $>2$ and the vast majority have size $1$ or $2$. Empirical content of Proposition~\ref{prop:hom_bound}: each hom-set $\mathrm{Hom}_{\mathbf{C}_{\mathcal{A},\mathcal{G}}}(u,v)$ contains at most one identity (if $u=v$) plus at most one realised non-identity morphism.
    (\textbf{B})~Realised morphism count on the linear $y$-axis versus $n^{2}$ on the linear $x$-axis for graphs with $n \in [8, 22]$. Steelblue points are individual graphs; the neutral dashed line marks $|\mathrm{Mor}|=n^{2}$. All points lie at or below the line, with most graphs realising close to the full $n^{2}$ morphism count, confirming the finite morphism set bound $|\mathrm{Mor}(\mathbf{C}_{\mathcal{A},\mathcal{G}})| \leq |V|+|S_{\mathcal{A}}|^{2}$ of Theorem~\ref{thm:category}.
    (\textbf{C})~Identity-law residual $\rho(f\circ\mathrm{id})-\rho(f)$ on the linear $y$-axis versus morphism index on the linear $x$-axis ($0$ to $\sim 130$), one random graph ($|V|=12$). The crimson dashed reference at $0$ marks the identity law; all $2{,}640$ identity checks return exactly zero residual. Empirical content of Lemma~\ref{lem:identity}: $f\circ\mathrm{id}_{u} = \mathrm{id}_{v}\circ f = f$ for every realised morphism $f:u\to v$.
    (\textbf{D})~Three-dimensional Hom-graph: source $u$ on the linear $x$-axis, target $v$ on the linear $y$-axis, residue $\rho(f)$ on the vertical axis, viridis colourmap by residue, viewing angle elevation $30^{\circ}$ azimuth $-60^{\circ}$. Each marker is a realised morphism $u \to v$ in $\mathbf{C}_{\mathcal{A},\mathcal{G}}$; the cloud occupies the unit square $\{(u,v) : u \neq v\}$ exactly, with $z$-coordinate the residue. Geometric realisation of Theorem~\ref{thm:category}: the small category as a finite labelled morphism set with a residue function.}
    \label{fig:panel-category}
\end{figure*}

% =====================================================================
\section{The Operational Structure Theorem}
\label{sec:main}
% =====================================================================

We now state and prove the main result: a finite agent engaging a
finite graph cannot avoid admitting an operational structure in the
following precise sense.

\subsection{Operational structure}

The previous section established three things: every realised
distinguishability act has residue $\geq \floor > 0$
(\cref{thm:floor}); residues compose subadditively
(\cref{thm:residue_additive}); the realised acts form a small
category $\mathbf{C}_{\Agent,\Graph}$ (\cref{thm:category}). The
operational structure packages these into a single object.

\begin{definition}[Residue equivalence and residue quotient]
\label{def:residue_equiv}
Two morphisms $f, g$ of $\mathbf{C}_{\Agent,\Graph}$ with
$\Res(f) = \Res(g)$ are \emph{residue-equivalent}, written
$f \equiv_\Res g$. The relation $\equiv_\Res$ is an equivalence on
each hom-set; the \emph{residue quotient}
$\widetilde{\mathbf{C}}_{\Agent,\Graph}$ identifies
residue-equivalent morphisms.
\end{definition}

\begin{theorem}[Residue cells are the natural unit of operation]
\label{thm:residue_cells}
The residue equivalence classes of $\equiv_\Res$ are
\emph{operation-cells}: regions of $\mathbf{C}_{\Agent,\Graph}$
within which morphisms are agent-indistinguishable. The residue
quotient $\widetilde{\mathbf{C}}_{\Agent,\Graph}$ is not a
remedial collapse of an antisymmetry failure but the natural
operational object: operations are individuated to their cell, not
to their label.
\end{theorem}

\begin{proof}
By \cref{thm:floor_from_infinitude}, individuation against the
non-completable whole proceeds by depositing residue
$\geq \floor > 0$; the agent's identification of a morphism with the
act it performs is to the resolution $\floor$ of its residue, not
finer. Two morphisms $f, g$ with $\Res(f) = \Res(g)$ deposit
indistinguishable measure differentials in the agent's record (the
agent's separator measure on $f$ equals that on $g$), so the
agent's finite descriptive scheme \ref{F:fin_descr} cannot
discriminate them by cost. The residue cells $[f]_{\equiv_\Res}$
are therefore the regions of operational indistinguishability under
the floor; the quotient
$\widetilde{\mathbf{C}}_{\Agent,\Graph}$ identifies them as the
natural unit of operation, and the residue preorder descends to a
genuine partial order on the cells. The cells are bounded below in
diameter by $\floor$ in the cost-metric on operations, mirroring
\cref{thm:floor_from_infinitude}. $\qed$
\end{proof}

\begin{remark}[Preorder collapse is cell-truth]
\label{rem:preorder_cell_truth}
The antisymmetry failure of $\preceq$ on $\mathbf{C}_{\Agent,
\Graph}$ is not a defect to be patched. It is the operational form
of cell-truth: the agent's record assigns a residue to each
operation; operations sharing a residue are the same operation to
the resolution of the floor; the natural object on which the order
lives is the residue-cell. To insist on antisymmetry over
$\mathbf{C}_{\Agent,\Graph}$ would be to demand point-truth on
operations, which \cref{thm:floor_from_infinitude} forbids. The
residue quotient is the operational structure's cell-structure, and
the partial order on the quotient is the unique order forced by
the conjunction of the three premises.
\end{remark}

\begin{figure*}[!htbp]
    \centering
    \includegraphics[width=\textwidth]{figures/panel_07_preorder.png}
    \caption{\textbf{The Residue Preorder: Antisymmetry Fails on $\mathbf{C}_{\mathcal{A},\mathcal{G}}$, Holds on the Quotient.}
    (\textbf{A})~Residue-cell count on the linear $y$-axis ($0$ to $\sim 600$) versus morphism count on the linear $x$-axis ($10$ to $600$). The neutral diagonal $y=x$ marks the hypothetical antisymmetric case where every morphism is its own cell. Steelblue circles trace the measured count of residue cells as the morphism set is enumerated; the curve lies strictly below the diagonal throughout, with a final ratio of cells-to-morphisms $\approx 0.5$. Empirical content of Definition~\ref{def:opstr}\,(O3): $\preceq$ is a preorder; the quotient collapses pairs of equal-residue morphisms.
    (\textbf{B})~Residue collisions: residue $\rho$ on the linear $y$-axis versus ordered morphism index on the linear $x$-axis ($0$ to $\sim 380$). Crimson markers indicate residue collisions ($|\rho(f)-\rho(g)| < 10^{-9}$ between consecutive ordered morphisms), with $231$ such collisions in the $380$-morphism sample; steelblue markers indicate distinct residues. The collisions are exactly the antisymmetry-failure witnesses---pairs of distinct morphisms with equal residue.
    (\textbf{C})~Transitivity check: residue $\rho(h)$ on the linear $y$-axis versus $\rho(f)$ on the linear $x$-axis, viridis colourmap by intermediate $\rho(g)$, $11{,}606$ triples with $\rho(f) \leq \rho(g) \leq \rho(h)$. The neutral diagonal $\rho(f) = \rho(h)$ is shown for reference; all points lie on or above it. Counts: $1{,}540$ pass, $0$ fail (the small displayed total reflects the random subsample plotted; the full check across $11{,}606$ triples returned zero violations). Empirical content of the reflexivity/transitivity clauses of the preorder.
    (\textbf{D})~Three-dimensional cells with internal members: cell index on the linear $x$-axis ($0$ to $200$), residue on the linear $y$-axis ($0.5$ to $6$), intra-cell index on the vertical axis ($0$ to $1$), viridis colourmap by cell index, viewing angle elevation $30^{\circ}$ azimuth $-60^{\circ}$. Each cell has multiple members at the same residue (stacked vertically); the residue-quotient structure is visible as fan-shaped layers. Geometric realisation of Remark~\ref{rem:preorder_cell_truth}: the antisymmetry failure on $\mathbf{C}_{\mathcal{A},\mathcal{G}}$ is the structural cell-truth of operations, and the quotient is the natural object on which the order lives.}
    \label{fig:panel-preorder}
\end{figure*}


\begin{definition}[Operational structure]
\label{def:opstr}
The \emph{operational structure} on a finite graph $\Graph$
relative to a finite agent $\Agent$ is the quintuple
\[
  \OpStr_\Agent \;=\;
  (\mathbf{C}_{\Agent,\Graph},\; \Res,\; \preceq,\; \floor,\;
  \mathrm{Realisable})
\]
where:
\begin{enumerate}[label=(O\arabic*),leftmargin=2.6em,nosep]
\item\label{O:cat} $\mathbf{C}_{\Agent,\Graph}$ is the small
category of \cref{def:dcat}.
\item\label{O:residue} $\Res: \mathrm{Mor}(\mathbf{C}_{\Agent,
\Graph}) \to \R_{\geq 0}$ is the residue function of
\cref{def:residue}, subadditive under composition
(\cref{thm:residue_additive}).
\item\label{O:preorder} $\preceq$ is the residue \emph{preorder}:
$f \preceq g$ iff $\Res(f) \leq \Res(g)$. On the residue quotient
$\widetilde{\mathbf{C}}_{\Agent,\Graph}$ this descends to a partial
order in the standard sense.
\item\label{O:floor} For every non-identity morphism $f$,
$\Res(f) \geq \floor > 0$, where $\floor$ is the graph-level floor
of \cref{thm:floor}.
\item\label{O:budget} $\mathrm{Realisable}: \mathrm{Mor}(\mathbf{C}_
{\Agent,\Graph}) \to \{0,1\}$ marks each morphism as realised or
not, subject to the budget constraint
$\sum_{f\text{ realised}} c(f) \leq C_\Agent$.
\end{enumerate}
\end{definition}

\begin{remark}[Why these five components]
\label{rem:five}
The category (\ref{O:cat}) captures composition; the residue
function (\ref{O:residue}) captures cost; the preorder
(\ref{O:preorder}) captures comparison; the floor (\ref{O:floor})
captures the irreducible minimum; the realisability marker
(\ref{O:budget}) captures the budget-constrained partiality. None
of these can be dropped without losing structural content forced by
the premises, and none requires more than the premises supply.
Together they describe what the agent can do, at what cost, in what
order, with what minimum, within what budget.
\end{remark}

\begin{theorem}[Operational Structure Theorem]
\label{thm:opstructure}
Under \cref{ax:graph} and \cref{ax:agent}, every distinguishability
relation $\sim$ that the agent realises on $\Graph$ determines a
unique operational structure $\OpStr_\Agent$. The structure is
forced by the conjunction of the two premises: each of its five
components is determined by $(\Agent, \Graph)$ together with the
agent's path-choices, and no component can be denied without
denying one of the premises.
\end{theorem}

\begin{proof}
We verify each clause and the uniqueness and forcing claims.

\emph{Category~\ref{O:cat}.} By \cref{thm:category},
$\mathbf{C}_{\Agent,\Graph}$ is a small category with object set
$\Vrt$, identity morphisms on each object, partial associative
composition. The morphisms include the identities and the
agent's realised distinguishability acts together with their paths.

\emph{Residue~\ref{O:residue}.} The residue function of
\cref{def:residue} assigns to each morphism the sum of edge
measures along its path. Subadditivity under composition is
\cref{thm:residue_additive}.

\emph{Preorder~\ref{O:preorder}.} The relation $f \preceq g$ iff
$\Res(f) \leq \Res(g)$ is reflexive ($\Res(f) \leq \Res(f)$) and
transitive ($\Res(f) \leq \Res(g) \leq \Res(h)$ implies
$\Res(f) \leq \Res(h)$). It is not antisymmetric on
$\mathbf{C}_{\Agent,\Graph}$ itself: distinct morphisms can have
equal residue, in which case both $f \preceq g$ and $g \preceq f$
hold without $f = g$. The preorder descends to a partial order on
the residue quotient
$\widetilde{\mathbf{C}}_{\Agent,\Graph} = \mathbf{C}_{\Agent,
\Graph}/\!\equiv_\Res$, where antisymmetry is recovered by
construction.

\emph{Floor~\ref{O:floor}.} By \cref{thm:floor}, the graph-level
floor $\floor = \min\{w(e) : e \in \Edg, w(e) > 0\}$ is a property
of $\Graph$ alone, and every realised non-identity morphism has
residue $\geq \floor$. Identity morphisms have residue $0$, so the
floor applies only to non-identity morphisms.

\emph{Realisability~\ref{O:budget}.} The realisability marker is
determined by the agent's actual run on $\Graph$, subject to
\cref{F:fin_cost}. Compositions of realised morphisms are
themselves realised only when the agent can record the composite
within the remaining budget; this is the partial-composition
phenomenon noted in \cref{rem:why_category}.

\emph{Uniqueness.} Each of the five components is determined by
$(\Agent, \Graph)$ together with the agent's path-choices: the
category is fixed by the realised acts; the residue function is
fixed by $w$; the preorder is fixed by residue comparison; the
floor is fixed by $\min\{w(e) : w(e) > 0\}$; the realisability
marker is fixed by the agent's run subject to the budget. Hence
$\OpStr_\Agent$ is unique given $(\Agent, \Graph)$ and the path
record.

\emph{Forcing.} The five clauses follow from the premises alone:
rejecting~\ref{O:cat} requires either infinite objects (denial of
\cref{ax:graph}) or non-associative composition (denial of
\cref{lem:assoc}, which uses only the agent's deterministic
recording, \cref{F:fin_descr});
rejecting~\ref{O:residue} requires negative or undefined edge
measures (denial of \cref{def:graph});
rejecting~\ref{O:preorder} requires denying the standard order on
$\R_{\geq 0}$;
rejecting~\ref{O:floor} requires zero-floor distinguishability acts
(denial of \cref{thm:floor}, hence of finiteness of $\Edg$);
rejecting~\ref{O:budget} requires unbounded budget (denial of
\cref{F:fin_cost}). The structure is forced. $\qed$
\end{proof}

\begin{remark}[Strengthening to a partial monoid]
\label{rem:partial_monoid}
If the agent's budget $C_\Agent$ is large enough that composition
of realised morphisms is always realised (no budget exhaustion on
composites), the category becomes a small category in which every
pair of composable morphisms composes---a one-object case of this
being a monoid. The general statement above does not assume this;
it admits the budget-exhausted case where some composites are not
realised. The strengthened result is a corollary under the
additional hypothesis $C_\Agent = \infty$ (or sufficient for
closure), which we do not adopt by default.
\end{remark}

\begin{remark}[Path-choice dependence]
\label{rem:path_choice}
$\OpStr_\Agent$ depends on the agent's path-choices because two
different paths between the same vertices can give different
residues. This is genuine: an agent choosing a longer path pays
more residue. We do not normalise to the shortest path; the
operational structure records what the agent actually did, not
what it could have done. To compare agents at the same graph, one
fixes a path-choice rule (e.g.~shortest-path policy) and the
remaining freedom collapses.
\end{remark}

\subsection{Consequences of the theorem}

\begin{corollary}[Operations are determined up to labelling and
path-choice]
\label{cor:not_chosen}
The category $\mathbf{C}_{\Agent,\Graph}$, the residue function
$\Res$, the preorder $\preceq$, the floor $\floor$, and the
realisability marker are determined by $(\Agent, \Graph)$ together
with the agent's path-choices. Choosing an operational structure
for a finite agent on a finite graph is choosing labels for
elements of a structure that already exists and, where the agent
has a choice of path between vertices, choosing the path; it is
not creating the structure.
\end{corollary}

\begin{proof}
Immediate from the uniqueness clause of \cref{thm:opstructure}.
The freedom is two-fold: labelling of objects and morphisms
(arbitrary names), and path-choice between vertices when multiple
paths exist (\cref{rem:path_choice}). The structure underneath is
fixed. $\qed$
\end{proof}

\begin{corollary}[Finite-system engineering is partly a discovery
task]
\label{cor:discovery}
Designing operations for a finite system engaged by finite agents
is partly the task of discovering the structure that
$(\Agent, \Graph)$ determines, and partly the task of fixing
path-choice and labelling conventions. Two designs disagreeing
about which operations to admit are disagreeing about one or more
of: which finite graph the system is modelled as; which finite
agent the program is; which path-choice policy governs route
selection.
\end{corollary}

\begin{proof}
By \cref{cor:not_chosen}, the structure is determined up to
labelling and path-choice. Differences between designs reduce to
differences in $(\Agent, \Graph)$ or in the path-choice policy.
Pure label differences do not affect the structure. $\qed$
\end{proof}

\begin{corollary}[Non-identity morphisms have residue
$\geq \floor$]
\label{cor:no_below}
Every non-identity morphism of $\mathbf{C}_{\Agent,\Graph}$ has
residue $\geq \floor > 0$. The floor is attained: morphisms of
residue exactly $\floor$ exist whenever some single-edge path of
the agent realises a distinction along an edge of weight $\floor$.
The floor cannot be crossed downward by any realisable composition.
\end{corollary}

\begin{proof}
The lower bound is~\ref{O:floor}. Attainment: pick any edge $e$ of
$\Graph$ with $w(e) = \floor$ (such an edge exists by the
definition of $\floor$ in \cref{thm:floor}). If the agent realises
the distinguishability across this single edge, the resulting
morphism has residue $\floor$. Non-crossable downward:
\cref{thm:residue_additive} gives subadditivity, so the residue of
a composite is bounded below by the sum of residues of its
non-identity components, each $\geq \floor$; the sum cannot be
less than $\floor$ unless all components are identities, in which
case the composite is itself the identity. $\qed$
\end{proof}


% =====================================================================
\section{Classification of Operational Structures}
\label{sec:classification}
% =====================================================================

We classify the operations arising in $\OpStr_\Agent$ by the
\emph{locus of their effect}. An operation acts on at most three
loci: the graph $\Graph$, the equivalence structure
$\sim$ on $\Vrt$, and the agent's internal state $S_\Agent$. We
classify by which loci an operation touches. This yields four
canonical types; we prove that these four exhaust the realisable
operations and give the residue lower bound for each.

\subsection{The four loci of effect}

\begin{definition}[Effect tuple]
\label{def:effect_tuple}
The \emph{effect tuple} of a realised operation $f$ is the triple
\[
  \mathrm{eff}(f) \;=\; (\Delta\Graph(f),\; \Delta{\sim}(f),\;
  \Delta S_\Agent(f))
\]
where $\Delta\Graph(f) \in \{0, 1\}$ records whether $f$ modifies
$\Graph$ (vertex/edge set or $w$); $\Delta{\sim}(f) \in \{0, 1\}$
records whether $f$ modifies the agent's equivalence relation on
$\Vrt$; and $\Delta S_\Agent(f) \in \{0, 1\}$ records whether $f$
modifies the agent's internal state $S_\Agent$ beyond what
$\Delta\Graph$ and $\Delta{\sim}$ already entail.
\end{definition}

\begin{remark}[Internal effects beyond the graph]
\label{rem:internal_beyond}
The third component $\Delta S_\Agent$ accounts for internal
bookkeeping not visible in $\Graph$ or $\sim$: counter increments,
timer updates, allocation of internal registers, logging of
realisations, garbage collection of stale records. By
\cref{F:fin_state} the internal state is finite; changes to it are
finite-step. By the finite-budget premise \cref{F:fin_cost} these
changes too consume part of $C_\Agent$, even when no
graph-side or equivalence-side effect occurs.
\end{remark}

\subsection{The four types}

\begin{definition}[Operation type]
\label{def:op_type}
A realised operation $f$ has type:
\begin{enumerate}[label=(T\arabic*),leftmargin=2.6em,nosep]
\item\label{T:alg} \emph{Algebraic} if
$\mathrm{eff}(f) = (0, 1, *)$: $f$ changes only the equivalence
relation on $\Vrt$, not the graph itself.
\item\label{T:dyn} \emph{Dynamical} if
$\mathrm{eff}(f) = (1, *, *)$: $f$ modifies $\Graph$.
\item\label{T:comp} \emph{Computational} if
$\mathrm{eff}(f) = (0, 0, *)$ and $f$ produces a binary output
(yes/no decision) by traversing a finite portion of $\Graph$
without modifying $\Graph$ or $\sim$.
\item\label{T:int} \emph{Internal} if
$\mathrm{eff}(f) = (0, 0, 1)$ and $f$ produces no externally-
visible output: $f$ updates only $S_\Agent$.
\end{enumerate}
The wildcard $*$ in \ref{T:alg} and \ref{T:dyn} indicates that
internal bookkeeping is permitted but not required.
\end{definition}

\paragraph{Algebraic.} Equivalence-class formation, vertex
relabelling under a graph automorphism, partition refinement.
Algebraic operations re-express the agent's distinguishability
relation without editing the graph or emitting output.

\paragraph{Dynamical.} Graph edits: adding or removing vertices,
adding or removing edges, modifying $w$. Dynamical operations
transition $\Graph_t \to \Graph_{t+1}$.

\paragraph{Computational.} Queries: membership tests, reachability
checks, path-existence checks, equivalence tests. A computational
operation returns a binary output by reading the graph.

\paragraph{Internal.} Pure internal-state operations: counter
increments, timer updates, allocation, logging, internal cache
invalidation, garbage collection of agent-side records. An internal
operation costs budget (it consumes finite agent resources) but
deposits no residue on $\Graph$ and produces no externally-visible
output.

\begin{lemma}[Type-specific residue bounds]
\label{lem:type_residues}
\begin{enumerate}[nosep,leftmargin=*]
\item Algebraic operations (\ref{T:alg}): residue $\geq \floor$
when the equivalence change requires reading at least one edge of
weight $\geq \floor$; residue $0$ when the equivalence change is a
relabelling using no edge information (e.g.~renaming with no
distinguishability assertion).
\item Dynamical operations (\ref{T:dyn}): residue $\geq \floor$
when the edit affects at least one edge of weight $\geq \floor$;
residue $0$ when the edit touches only null edges (a vacuous case
that cannot contribute to any distinguishability,
\cref{thm:floor}).
\item Computational operations (\ref{T:comp}): residue $\geq \floor$
when the query reads at least one positive-measure edge; residue
$0$ for queries answered without reading any positive-measure edge
(e.g.~constant queries, queries decidable from $S_\Agent$ alone).
\item Internal operations (\ref{T:int}): residue $0$ on $\Graph$
(by definition no graph access), but cost $\geq c_0 > 0$ on the
agent's budget by \cref{F:fin_cost}.
\end{enumerate}
\end{lemma}

\begin{proof}
Each clause is immediate from the definitions and from
\cref{thm:floor}. The residue bound applies to graph-side
contributions; budget cost is bounded below by $c_0$ for any
recorded operation, including internal ones. $\qed$
\end{proof}

\subsection{Four-type exhaustiveness}

\begin{theorem}[Four-type exhaustiveness]
\label{thm:three_types}
Every realised operation $f \in \mathrm{Mor}(\mathbf{C}_{\Agent,
\Graph})$ has effect tuple in
$\{0,1\}^3 \setminus \{(0,0,0)\}$ and is therefore one of the four
types of \cref{def:op_type}, or a composite whose components fall
into these types. The all-zero tuple $(0,0,0)$ corresponds to the
identity morphism, which is not a non-trivial operation.
\end{theorem}

\begin{proof}
The effect tuple has $2^3 = 8$ possible values. We classify each:
\begin{enumerate}[nosep,leftmargin=*]
\item $(0,0,0)$: no effect anywhere; this is the identity morphism
$\mathrm{id}$.
\item $(0,0,1)$: only internal state changes; type \ref{T:int}
(internal).
\item $(0,1,0)$ and $(0,1,1)$: equivalence relation changes
without graph modification; type \ref{T:alg} (algebraic). The
internal-state component is permitted as bookkeeping
(\cref{rem:internal_beyond}).
\item $(1,0,0)$, $(1,0,1)$, $(1,1,0)$, $(1,1,1)$: graph
modifications, possibly accompanied by equivalence or internal
changes; type \ref{T:dyn} (dynamical), with the wildcard absorbing
the other components.
\end{enumerate}
The eighth case, computational output, is not encoded in
$\mathrm{eff}(f)$ as a fourth component because output by itself
is an internal-state effect (the agent's internal state records the
output). A computational operation that reads $\Graph$ without
modifying it and writes its output to $S_\Agent$ has effect
$(0,0,1)$, which is also the internal type; the distinguishing
feature of computational versus internal is whether the operation
\emph{reads the graph} during its execution. We refine: type
\ref{T:comp} (computational) is the subtype of $(0,0,1)$ in which
the agent traverses at least one edge of $\Graph$ during the
operation; type \ref{T:int} (internal) is the subtype of $(0,0,1)$
in which no graph edge is traversed.

This gives four mutually exclusive types covering all
non-identity effect tuples. Any realised operation is either an
identity, falls into one of the four types, or is a composite
whose components do. $\qed$
\end{proof}

\begin{theorem}[Four-type taxonomy is the scale homomorphism]
\label{thm:scale_homomorphism}
The four types of \cref{def:op_type} are not an ad hoc taxonomy
but the image of a structural homomorphism between the
operations-scale and the contact-graph scale that the agent
engages. The four loci of effect---graph $\Graph$, equivalence
relation $\sim$, internal state $S_\Agent$, traversed-edges
trace---are the agent-scale projections of the four-place
structure $(\text{unit},\,\text{whole},\,\text{selector},\,\text{
invariant})$ at the operations scale. Specifically, the assignment
\[
  \begin{array}{rcl}
  \text{whole at op-scale} & \longmapsto & \Graph \\
  \text{individuating selector} & \longmapsto & \sim \\
  \text{invariant residue} & \longmapsto & S_\Agent \\
  \text{committed history} & \longmapsto & \text{edges traversed}
  \end{array}
\]
maps each role of the operations-scale's four-place structure to
the locus of effect that records it on the agent-scale, and the
four operation types are exactly the four singletons (and their
compositions) in this image.
\end{theorem}

\begin{proof}
At the operations scale, the agent's individuating activity over
$\Graph$ admits a four-place description: a whole the agent
engages, a selector that picks out the next individuation, an
invariant the selector deposits, and a committed history of
selections made. These are the four roles any finite individuating
agent in a non-completable whole exhibits
(\cref{ax:noncompletable}, \cref{ax:agent}), and they are the
operational counterpart of the four-place structure
$(\text{part},\,\text{whole},\,\text{gate},\,\text{invariant})$
that the contact-graph theory uses across scales.

Each role surfaces on the agent-scale at exactly one locus of
effect. The \emph{whole} surfaces as modifications of $\Graph$
itself (dynamical, \ref{T:dyn}): the only way an operation engages
the whole is to change the graph it models. The \emph{selector}
surfaces as modifications of $\sim$ (algebraic, \ref{T:alg}): the
selector is the equivalence relation through which the agent
individuates parts; changing it is changing what is treated as one.
The \emph{invariant residue} surfaces as modifications of
$S_\Agent$ (internal, \ref{T:int}): the invariant the agent
deposits is recorded internally. The \emph{committed history}
surfaces as traversal of graph edges (computational, \ref{T:comp}):
querying produces a record of edges read, the operation's commitment.

The four types of \cref{def:op_type} are therefore the four
singletons in the image of the role-to-locus assignment, and any
realised operation is a composite of these (the agent acts on at
most all four loci simultaneously). The taxonomy is forced by the
scale homomorphism; it is not a definition with attached
completeness claim but a theorem about which loci the four roles
surface at. $\qed$
\end{proof}

\begin{remark}[Why the taxonomy is structural, not stipulated]
\label{rem:taxonomy_honest}
\Cref{thm:scale_homomorphism} converts the four-type taxonomy from
a stipulation about loci of effect into a structural fact about
roles at one scale projecting to loci at the next. The same
four-place structure recurs at every scale at which an agent
engages a non-completable whole: a unit individuating against a
whole, with a selector and an invariant, over a committed history.
At each scale the four roles surface as the four loci, and the
operations are typed by which roles they exercise. The
exhaustiveness is structural---there are exactly four roles---not
contingent on the locus model. An extension to additional roles
(e.g.\ inter-agent communication as a fifth role) would extend the
taxonomy by exactly one type, but the conjunction of the three
premises supplies four roles, hence four types.
\end{remark}

\begin{remark}[Composition across types]
\label{rem:cross_type}
Operations of different types compose where defined: a query
followed by an edit followed by a relabelling is a single
composite morphism in $\mathbf{C}_{\Agent,\Graph}$. The
classification labels the type of each component; composition
operates on the composites. As before, composition is partial
(\cref{rem:why_category}); not every composite of realised
operations is itself realised.
\end{remark}

\begin{figure*}[!htbp]
    \centering
    \includegraphics[width=\textwidth]{figures/panel_08_scale_homomorphism.png}
    \caption{\textbf{Scale Homomorphism: Four Roles Surface as Four Loci.}
    (\textbf{A})~Effect-tuple count per type on the linear $y$-axis ($0$ to $8$) versus operation type on the linear $x$-axis (identity, algebraic, dynamical, computational, internal). Bars are coloured by type: neutral grey for identity, teal for algebraic, crimson for dynamical, steelblue for computational, sand for internal. Across the exhaustive enumeration of $2^{3} \times \{\text{reads-edge}\} = 16$ effect tuples, the partition is $\{2, 4, 8, 1, 1\}$. Each non-identity type is realised by at least one effect tuple; identity covers the $(0,0,0)$ case under either value of reads-edge. Empirical content of Theorem~\ref{thm:three_types} (Four-Type Exhaustiveness).
    (\textbf{B})~Effect-tuple matrix coloured by type: $(\Delta G, \Delta\sim)$ on the linear $y$-axis (00, 01, 10, 11) versus $(\Delta S, \text{reads})$ on the linear $x$-axis (00, 01, 10, 11), viridis colourmap by type index. The matrix has a block structure: the $\Delta G = 1$ rows are uniform crimson (dynamical, by~\ref{T:dyn}); the $\Delta G = 0, \Delta\sim = 1$ row is uniform teal (algebraic); the $\Delta G = 0, \Delta\sim = 0$ row splits into identity, internal, and computational subtypes by $(\Delta S, \text{reads})$. Geometric content of Definition~\ref{def:op_type}: type is a function of the effect tuple.
    (\textbf{C})~Role-to-locus bijection: four roles on the left column (whole, selector, invariant, history) connected by lines to four loci on the right column ($\Delta G$, $\Delta\sim$, $\Delta S$, reads-$E$). Each role maps to exactly one locus; the assignment is a bijection. Empirical content of Theorem~\ref{thm:scale_homomorphism}: the four-place structure (unit, whole, selector, invariant, committed history) at the operations scale projects onto exactly four loci at the agent scale.
    (\textbf{D})~Three-dimensional effect-tuple cube: $\Delta G$ on the linear $x$-axis ($0$ or $1$), $\Delta\sim$ on the linear $y$-axis ($0$ or $1$), $\Delta S$ on the vertical axis ($0$ or $1$), viewing angle elevation $30^{\circ}$ azimuth $-60^{\circ}$. Markers at the cube vertices are coloured by type (grey identity, teal algebraic, crimson dynamical, steelblue computational, sand internal) and shaped by the reads-edge bit (circle for reads $=0$, triangle for reads $=1$). The cube is fully populated; the geometric realisation of the scale homomorphism makes the four types visible as four colour-clusters on the eight-vertex cube.}
    \label{fig:panel-scale-homomorphism}
\end{figure*}




% =====================================================================
\section{Composition Inflation: Trajectory Counting in the
Operational Structure}
\label{sec:inflation}
% =====================================================================

The Operational Structure Theorem (\cref{thm:opstructure}) bounds the
size of $\mathbf{C}_{\Agent,\Graph}$ at $|\Vrt| + |S_\Agent|^2$
morphisms. This is the bound on \emph{individual} realised acts. The
agent, however, does not realise acts in isolation: it commits a
sequence of acts, and the sequence carries its own categorical
structure. We now derive the count of distinguishable operation
\emph{trajectories} as a function of trajectory length, and show
that the count grows exponentially against a linearly-growing
committed cost.

The result is a closed-form formula
$T(n, d) = d \cdot (d+1)^{n-1}$ for the number of categorically
distinguishable operation trajectories of length $n$ in a structure
with $d$ operation types. It is proved here from first principles
using only the four-type taxonomy
(\cref{thm:scale_homomorphism}) and the agent's commitment record
(\cref{F:fin_descr}, \cref{F:fin_cost}). The closed form supplies
the precise rate at which the operational structure's trajectory
space inflates relative to the agent's committed-act count.

\subsection{Operation trajectories and commitment blocks}

\begin{definition}[Operation trajectory]
\label{def:trajectory}
An \emph{operation trajectory} of length $n$ in $\OpStr_\Agent$ is
an ordered sequence
\[
  \tau \;=\; (f_1, f_2, \ldots, f_n)
\]
of $n$ realised non-identity morphisms $f_i \in \mathbf{C}_{\Agent,
\Graph}$, recorded by the agent's descriptive scheme
(\cref{F:fin_descr}) in commitment order. Each $f_i$ carries a type
$t(f_i) \in \{\text{algebraic}, \text{dynamical}, \text{computational},
\text{internal}\}$ from \cref{def:op_type}.
\end{definition}

\begin{definition}[Commitment block; type-segmentation]
\label{def:block}
A \emph{commitment block} of a trajectory $\tau$ is a maximal
contiguous sub-sequence of acts that the agent records as a single
recorded unit. Formally, given $\tau = (f_1, \ldots, f_n)$, a
\emph{block structure} on $\tau$ is a composition
$(c_1, c_2, \ldots, c_k)$ of $n$ with $c_1 + \cdots + c_k = n$,
$c_i \geq 1$, indicating that the first $c_1$ acts form the first
block, the next $c_2$ acts form the second block, and so on. The
\emph{type-labeling} of a block structure assigns each block a type
$t_j \in \{1, 2, 3, 4\}$ (one of the four operation types of
\cref{def:op_type}). The number of blocks $k$ is the
\emph{commitment depth} of the trajectory.
\end{definition}

\begin{remark}[Why block structure is data, not derived]
\label{rem:why_blocks}
A naive view would treat the type sequence $(t(f_1), \ldots, t(f_n))$
as the full trajectory data and read block boundaries off as
positions of type change. This conflates two distinct
trajectories: the agent committing acts in pairs $(f_1, f_2)$ then
$(f_3, f_4)$ versus committing four separate acts $f_1, f_2, f_3,
f_4$ of the same type sequence. The two have different commitment
records: the first records two checkpoints; the second records
four. By \cref{F:fin_descr}, the agent's descriptive scheme
distinguishes these, so they are distinct trajectories in the
operational structure. The block structure is therefore part of the
trajectory data, not a derived quantity.
\end{remark}

\subsection{The composition count}

\begin{lemma}[Compositions of $n$]
\label{lem:compositions}
The number of compositions of a positive integer $n$ into $k$
positive parts is $\binom{n-1}{k-1}$. Summing over $k$, the total
number of compositions of $n$ is $2^{n-1}$.
\end{lemma}

\begin{proof}
A composition of $n$ into $k$ parts corresponds bijectively to a
choice of $k - 1$ dividers placed into the $n - 1$ gaps between
$n$ ordered unit objects: each divider marks a boundary between
adjacent parts. The number of such choices is $\binom{n-1}{k-1}$.
Summing,
\[
  \sum_{k=1}^{n} \binom{n-1}{k-1}
  \;=\; \sum_{j=0}^{n-1} \binom{n-1}{j}
  \;=\; (1+1)^{n-1}
  \;=\; 2^{n-1}
\]
by the binomial theorem applied to $(1+1)^{n-1}$. $\qed$
\end{proof}

\subsection{The composition-inflation theorem}

\begin{theorem}[Composition Inflation in Operation Space]
\label{thm:composition_inflation}
Let $\OpStr_\Agent$ have $d$ operation types (with $d = 4$ in the
canonical setting of \cref{def:op_type}). The number of
categorically distinguishable operation trajectories of length $n$
in $\OpStr_\Agent$ is
\[
  \boxed{\;T(n, d) \;=\; d \cdot (d + 1)^{n - 1}.\;}
\]
\end{theorem}

\begin{proof}
A trajectory of length $n$ is uniquely specified by:
\begin{enumerate}[nosep,leftmargin=*]
\item A block structure: a composition $(c_1, \ldots, c_k)$ of $n$
into $k$ parts, for some $1 \leq k \leq n$
(\cref{def:block}).
\item A type-labeling: an assignment of one of $d$ types to each
of the $k$ blocks.
\end{enumerate}
By \cref{lem:compositions}, the number of compositions of $n$ into
$k$ parts is $\binom{n-1}{k-1}$. For each such composition, the
$k$ type labels are chosen independently from $d$ options, yielding
$d^k$ labelings. Hence the number of trajectories with commitment
depth exactly $k$ is
\[
  N_k(n, d) \;=\; \binom{n-1}{k-1} \cdot d^k.
\]
Summing over commitment depth:
\begin{align*}
  T(n, d) &\;=\; \sum_{k=1}^{n} \binom{n-1}{k-1} \cdot d^k \\
  &\;=\; d \sum_{k=1}^{n} \binom{n-1}{k-1} \cdot d^{k-1} \\
  &\;=\; d \sum_{j=0}^{n-1} \binom{n-1}{j} \cdot d^{j}
  \quad [\text{substitute } j = k - 1] \\
  &\;=\; d \cdot (1 + d)^{n-1}
\end{align*}
by the binomial theorem applied to $(1 + d)^{n-1}$. $\qed$
\end{proof}

\begin{corollary}[Recurrence and initial condition]
\label{cor:inflation_recurrence}
$T(n, d)$ satisfies the recurrence
$T(n+1, d) = (d + 1) \cdot T(n, d)$ with initial condition
$T(1, d) = d$. Equivalently, each additional committed act
multiplies the trajectory count by $d + 1$.
\end{corollary}

\begin{proof}
$T(1, d) = d \cdot (d+1)^0 = d$, agreeing with the direct count
(one act, $d$ type choices). For $n \geq 1$,
$T(n + 1, d) = d \cdot (d+1)^{n} = (d+1) \cdot T(n, d)$. $\qed$
\end{proof}

\begin{remark}[Where the extra $+1$ comes from]
\label{rem:plus_one}
The exponential base $d + 1$ has a structural meaning. At each of
the $n - 1$ internal boundaries between consecutive committed acts,
the agent has two independent decisions: (i) start a new block or
continue the current one (binary, contributing the composition
factor); (ii) if starting a new block, choose its type (one of $d$).
Together these are $1 + d$ effective branches per boundary
($1$ for ``continue'', $d$ for ``branch into one of $d$ types''),
yielding base $d + 1$. This is the agent-scale reading of the
combinatorial constant.
\end{remark}

\subsection{Linear cost, exponential content}

\begin{theorem}[Asymmetry between cost and content]
\label{thm:cost_content_asymmetry}
Let $\Agent$ commit $n$ realised non-identity acts in $\OpStr_\Agent$.
Then:
\begin{enumerate}[label=\textup{(\roman*)},nosep,leftmargin=2.4em]
\item The cumulative residue (committed cost) is bounded below by
$n \cdot \floor > 0$ (\cref{cor:res_accumulate}); it grows
\emph{linearly} in $n$.
\item The number of categorically distinguishable operation
trajectories the commitment can realise is $T(n, d)
= d \cdot (d + 1)^{n - 1}$ (\cref{thm:composition_inflation});
it grows \emph{exponentially} in $n$ with base $d + 1$.
\end{enumerate}
The categorical content per unit committed cost is therefore
\[
  \frac{T(n, d)}{n \cdot \floor}
  \;=\; \frac{d \cdot (d + 1)^{n - 1}}{n \cdot \floor}
  \;\xrightarrow[n \to \infty]{}\; \infty,
\]
strictly superlinear in $n$.
\end{theorem}

\begin{proof}
(i) Each of the $n$ realised non-identity acts deposits residue
$\geq \floor$ by \cref{thm:floor}; summing gives $\sum_{i=1}^{n}
\rho(f_i) \geq n \floor$. This is the linear lower bound.
(ii) Immediate from \cref{thm:composition_inflation}.
The ratio $T(n, d) / (n \floor)$ has numerator growing as
$(d+1)^{n-1}$ and denominator growing as $n$; the exponential
dominates, so the ratio diverges. $\qed$
\end{proof}

\begin{corollary}[Logarithmic termination for content-bounded
processes]
\label{cor:log_termination}
Let a process running on $\OpStr_\Agent$ terminate when its
committed history realises $K$ distinguishable operation
trajectories. The number of committed acts required for
termination is
\[
  n_{\mathrm{term}}(K)
  \;=\; 1 + \left\lceil
    \log_{d+1}\!\left(\frac{K}{d}\right)
  \right\rceil,
\]
and the committed cost at termination is
\[
  \rho_{\mathrm{term}}(K)
  \;\geq\; n_{\mathrm{term}}(K) \cdot \floor
  \;=\; O(\log K).
\]
\end{corollary}

\begin{proof}
By \cref{thm:composition_inflation}, $K$ distinguishable
trajectories are realised at the smallest $n$ with
$T(n, d) \geq K$, i.e.~$d \cdot (d+1)^{n-1} \geq K$, equivalently
$(d+1)^{n-1} \geq K/d$, equivalently
$n \geq 1 + \log_{d+1}(K/d)$. The smallest integer satisfying this
is $n_{\mathrm{term}}(K) = 1 + \lceil \log_{d+1}(K/d) \rceil$. The
committed cost lower bound follows from
\cref{cor:res_accumulate}. $\qed$
\end{proof}

\begin{remark}[The termination claim in plain language]
\label{rem:termination_plain}
A process that needs ``enough categorical content to settle a
question'' terminates in $O(\log K)$ committed acts when the
content threshold is $K$. The agent does not need to commit more
acts to produce more trajectories; it commits the same act and
composition-inflation supplies the trajectories. The four-type
taxonomy (\cref{thm:scale_homomorphism}) is what makes the
exponential base $d + 1 = 5$ rather than $2$; the four loci of
effect multiply the agent's per-act categorical bandwidth.
\end{remark}

\begin{remark}[What this does and does not say]
\label{rem:inflation_scope}
\Cref{thm:composition_inflation} concerns the
\emph{categorical-content} measure of the operational structure: the
count of distinguishable operation trajectories. It does not say
any individual act is faster, cheaper, or more powerful---each act
still pays residue $\geq \floor$ (\cref{thm:floor}). It does not
say the cumulative cost grows logarithmically: the cumulative
residue is at least $n \floor$, which is linear in $n$. The
asymmetry is between two measures, both true: the cost-measure is
linear, the content-measure is exponential. Processes whose
termination depends on the content-measure (``produce $K$
trajectories'') terminate logarithmically; processes whose
termination depends on the cost-measure (``deposit total residue
$\geq R$'') terminate linearly. The framework is silent about which
measure governs any given process; that is a property of the
process specification.
\end{remark}

\subsection{Worked example: small cases}

\begin{example}
\label{ex:small_cases}
Direct enumeration for $d = 4$ verifies the formula on small
cases:
\begin{itemize}[nosep,leftmargin=*]
\item $n = 1$: $T(1, 4) = 4$. One act, four type choices.
Direct count: $\{(t) : t \in \{1,2,3,4\}\}$, $|\cdot| = 4$.
\item $n = 2$: $T(2, 4) = 4 \cdot 5 = 20$. Compositions of $2$ are
$(2)$ and $(1, 1)$. For composition $(2)$: $4$ labelings (one
block, one type each). For composition $(1, 1)$: $4 \times 4 =
16$ labelings (two independent block types). Total $4 + 16 = 20$.
\item $n = 3$: $T(3, 4) = 4 \cdot 25 = 100$. Compositions of $3$
are $(3)$, $(2, 1)$, $(1, 2)$, $(1, 1, 1)$; counts $4 + 16 + 16 +
64 = 100$.
\end{itemize}
\end{example}

\begin{example}[Scaling]
\label{ex:scaling}
For $d = 4$ and $n = 10$, the trajectory count is
$T(10, 4) = 4 \cdot 5^9 = 7\,812\,500$. A process that requires
$10^6$ distinguishable trajectories terminates at
$n_{\mathrm{term}}(10^6) = 1 + \lceil \log_5(2.5 \times 10^5)
\rceil = 1 + \lceil 7.72 \rceil = 9$ committed acts. The
committed cost lower bound is $9 \floor$---linear in $9$,
logarithmic in $10^6$.
\end{example}

% =====================================================================
\section{The Forced Structure in Detail}
\label{sec:detail}
% =====================================================================

We restate the main result as a structural recipe, making explicit
what $\OpStr_\Agent$ looks like once $(\Agent, \Graph)$ is fixed.

\subsection{The structure recipe}

Given $(\Agent, \Graph)$ satisfying \cref{ax:agent} and
\cref{ax:graph}, together with a path-choice policy for the agent:
\begin{enumerate}[label=(R\arabic*),leftmargin=2.6em,nosep]
\item\label{R:floor} Compute $\floor = \min \{w(e) : e \in \Edg,\,
w(e) > 0\}$. This is the graph-level floor.
\item\label{R:cat} Enumerate the agent's realised
distinguishability acts on $\Graph$; together with the identity on
each vertex, these are the morphisms of
$\mathbf{C}_{\Agent,\Graph}$.
\item\label{R:residue} For each morphism $f$, compute the residue
$\Res(f)$ as the path-measure sum.
\item\label{R:preorder} Order the morphisms by residue:
$f \preceq g$ iff $\Res(f) \leq \Res(g)$. Quotient by
residue-equivalence $\equiv_\Res$ to obtain a partial order on
$\widetilde{\mathbf{C}}_{\Agent,\Graph}$.
\item\label{R:classify} Classify each morphism by effect tuple
(\cref{def:effect_tuple}) into one of the four types
\ref{T:alg}--\ref{T:int}.
\item\label{R:budget} Mark each morphism as realised or not,
subject to the budget constraint $\sum c(f) \leq C_\Agent$.
\end{enumerate}

\begin{proposition}[The recipe is well-defined and terminates]
\label{prop:recipe}
The recipe \ref{R:floor}--\ref{R:budget} is well-defined and
terminates in finite time for any $(\Agent, \Graph)$ satisfying
the two premises.
\end{proposition}

\begin{proof}
\ref{R:floor} is a minimum over a finite non-empty set of positive
reals: well-defined, computable in $O(|\Edg|)$ time.
\ref{R:cat} enumerates a finite morphism set bounded by
$|\Vrt| + |S_\Agent|^2$, terminating in finite time.
\ref{R:residue} sums finitely many non-negative reals per morphism.
\ref{R:preorder} is a comparison on a finite totally ordered set
($\R_{\geq 0}$); residue-equivalence quotienting is a finite
partition.
\ref{R:classify} is a case distinction over four cases per
morphism.
\ref{R:budget} is a finite sum compared to a finite bound.
Each step terminates; the recipe terminates. $\qed$
\end{proof}

\subsection{What the structure looks like}

\begin{proposition}[Structural shape]
\label{prop:shape}
The operational structure $\OpStr_\Agent$ has the following shape:
\begin{enumerate}[nosep,leftmargin=*]
\item It is a small category (\cref{thm:category}) with finite
object set and finite morphism set, and with partial composition
constrained by the agent's budget.
\item Its non-identity morphisms have residues bounded below by
$\floor > 0$ (\cref{thm:floor}).
\item It carries a residue preorder, with identity morphisms at
residue $0$; on the residue quotient
$\widetilde{\mathbf{C}}_{\Agent,\Graph}$ the preorder descends to a
partial order.
\item Its morphisms decompose into four canonical types
(\cref{thm:three_types}): algebraic, dynamical, computational,
internal.
\end{enumerate}
\end{proposition}

\begin{proof}
Each clause is a restatement of an earlier result. The shape is
the conjunction of those results. $\qed$
\end{proof}

\subsection{Why the structure is operational rather than only
combinatorial}

The structure $\OpStr_\Agent$ is called \emph{operational} because
each morphism represents an act the agent can perform, with a
quantifiable cost (residue) and a quantifiable order (residue
preorder). It differs from a purely combinatorial structure---a
small category considered abstractly---in two ways:
\begin{enumerate}[nosep,leftmargin=*]
\item Each element is associated with a concrete act on the graph,
not merely an abstract symbol.
\item Composition is realised, not merely defined: composing two
acts produces a third act the agent can perform, with predictable
residue deposited on the graph.
\end{enumerate}
The operational character is essential to the result: we are not
classifying the agent's grammar but its realisable acts on a
concrete finite system. The two coincide only because the premises
force the agent's grammar to track its acts; in richer settings
the two diverge.

\begin{figure*}[!htbp]
    \centering
    \includegraphics[width=\textwidth]{figures/panel_05_partiality.png}
    \caption{\textbf{Partiality of Composition is Forced by Non-Completability, Not by Budget Exhaustion.}
    (\textbf{A})~Composable pair count on the linear $y$-axis (up to $\sim 3{,}000$) versus refinement stage $n$ on the linear $x$-axis ($2$ to $8$). Steelblue bars: composable pairs $(a,b),(b,c)$ for which the composite $(a,c)$ is realised. Crimson bars: composable pairs for which the composite is \emph{not} realised. At every stage tested, the unrealised count is substantial; the gap is not asymptotically closed by increasing $n$. Empirical content of Theorem~\ref{thm:partiality_forced}: composition in $\mathbf{C}_{\mathcal{A},\mathcal{G}}$ is genuinely partial.
    (\textbf{B})~Budget spent fraction on a log $y$-axis ($10^{-17}$ to $10^{1}$) versus stage $n$ on the linear $x$-axis ($2$ to $8$). Steelblue circles trace the agent's spent fraction; the crimson dashed reference at $1.0$ marks budget exhaustion. The measured spent fraction at every stage is $\sim 10^{-16}$ of budget---essentially nothing. Partiality is therefore \emph{not} an artefact of budget exhaustion; the agent's budget is virtually unconsumed at every stage. This is the decisive empirical content distinguishing Theorem~\ref{thm:partiality_forced} from the budget-exhaustion contingency.
    (\textbf{C})~Unrealised composable fraction on the linear $y$-axis ($0$ to $1$) versus stage $n$ on the linear $x$-axis ($2$ to $8$). Steelblue circles cluster around $0.5$; the neutral horizontal reference at $0.5$ is included for orientation. Roughly half of all composable pairs at every stage have unrealised composites, despite the budget being effectively unused---structural partiality.
    (\textbf{D})~Three-dimensional composability gap by stage: stage $n$ on the linear $x$-axis ($2$ to $8$), composite status (realised at $y=0$, unrealised at $y=1$) on the linear $y$-axis, count on the vertical axis. Steelblue ribbons mark realised composites; crimson ribbons mark unrealised composites. Both grow with $n$ in lockstep, confirming that partiality scales with the size of the agent's record without ever closing. Geometric realisation of Theorem~\ref{thm:partiality_forced}: the gap between the realised category and its composition closure is a non-empty sub-lattice at every finite stage.}
    \label{fig:panel-partiality}
\end{figure*}



% =====================================================================
\section{Falsifiability}
\label{sec:falsifiability}
% =====================================================================

\subsection{Where the argument can fail}

The main theorem (\cref{thm:opstructure}) is a conditional claim
under \cref{ax:graph} and \cref{ax:agent}. It can be falsified only
by:
\begin{enumerate}[nosep,leftmargin=*]
\item Showing the two premises jointly imply a contradiction (they
do not, by inspection: a finite graph and a finite agent obviously
coexist; we are reading this on one).
\item Showing some step in the derivation does not follow from the
premises (a mathematical refutation of one of the proved lemmas).
\item Showing the premises do not apply to a particular case (which
is honest: the framework declines to address that case).
\end{enumerate}

The substantive falsification target is (2): exhibiting a finite
agent realising distinguishability with zero floor on a finite
graph. We close this case in the next subsection.

\subsection{Zero-floor distinguishability is structurally impossible}

\begin{theorem}[Zero-floor impossibility]
\label{thm:zero_floor}
Under \cref{ax:graph} and \cref{ax:agent}, there is no agent--graph
pair $(\Agent, \Graph)$ admitting a realised distinguishability
act of residue $< \floor$, where $\floor = \min\{w(e) : e \in \Edg,
w(e) > 0\}$ is the graph-level floor of \cref{thm:floor}.
\end{theorem}

\begin{proof}
Suppose for contradiction such a pair exists, with a realised act
$u \not\sim v$ of residue $\Res(u, v) < \floor$. The residue is the
sum of edge measures along the agent's path. Each edge of the path
has measure $0$ or $\geq \floor$ (by the definition of $\floor$ as
the minimum positive edge measure in $\Graph$). Hence the sum is
$0$ or $\geq \floor$.

If the sum is $0$, every edge on the path has measure $0$. The
agent's descriptive scheme has no measure differential to record
the distinction (\cref{F:fin_descr}, see the realised-path argument
in the proof of \cref{thm:floor}), so the act is not realised in
the agent's record. The assumed realised act does not exist.

If the sum is $\geq \floor$, the residue is $\geq \floor$,
contradicting the assumption $\Res(u, v) < \floor$.

Either case contradicts the assumption. $\qed$
\end{proof}

\begin{corollary}[The main theorem is irrefutable from within the
premises]
\label{cor:irrefutable}
\Cref{thm:opstructure} cannot be refuted by an example satisfying
\cref{ax:graph} and \cref{ax:agent}. Refutation requires either
denying one of the premises or finding an error in a constituent
lemma; no example produced under the premises is a counterexample.
\end{corollary}

\begin{proof}
Refutation by example requires a $(\Agent, \Graph)$ satisfying the
premises in which the theorem fails. By the construction of
$\OpStr_\Agent$ from $(\Agent, \Graph)$ in the theorem's proof, and
the zero-floor impossibility (\cref{thm:zero_floor}), no such
example exists. $\qed$
\end{proof}

\subsection{Honest scope}

\begin{remark}[The argument is silent outside its premises]
\label{rem:silent}
We make no claim about agents that fail \cref{ax:agent} (putative
infinite agents, oracles, perfect observers), systems that fail
\cref{ax:graph} (continuous fields without finite discretisation,
infinite combinatorial structures), or wholes that fail
\cref{ax:noncompletable} (completable finite wholes that admit a
terminal stage). The argument has nothing to say about those cases.
This silence is intended: the framework is deliberately narrow,
sharp where the premises hold, mute elsewhere. This is the right
shape for a conditional architectural claim. It allows the result
to be definite without overreaching.
\end{remark}

\subsection{The Master Theorem}

We now show that the three premises are not merely sufficient for
the operational structure; they are equivalent to it, jointly. The
operational structure exists if and only if the conjunction of the
three premises holds; failure of any one premise dissolves at
least one structural feature.

\begin{theorem}[Operational Structure Master Theorem]
\label{thm:master}
Let an agent $\Agent$ engage a system modelled by a graph $\Graph$
within an ambient whole $\mathcal{M}$. The following are
equivalent:
\begin{enumerate}[label=\textup{(\roman*)},leftmargin=2.2em]
\item There exists a realisable distinguishability act
$u \not\sim v$ by $\Agent$ on $\Graph$ within $\mathcal{M}$.
\item \Cref{ax:graph}, \cref{ax:agent}, and
\cref{ax:noncompletable} all hold.
\item The full operational structure $\OpStr_\Agent$ of
\cref{def:opstr} is forced: a small category
$\mathbf{C}_{\Agent,\Graph}$ with strictly positive floor $\floor$
on non-identity residues, partial composition (forced by
\cref{thm:partiality_forced}), residue cells as the natural unit
(\cref{thm:residue_cells}), and the four-type taxonomy as scale
homomorphism (\cref{thm:scale_homomorphism}).
\end{enumerate}
\end{theorem}

\begin{proof}
$(\mathrm{i}) \Rightarrow (\mathrm{ii})$: A realised act
$u \not\sim v$ requires $u$ and $v$ to be distinguishable vertices
of $\Graph$, hence $\Graph$ must be a graph (\cref{ax:graph}) and
$\Agent$ must record the distinction (\cref{ax:agent}). The act's
residue is bounded below by \cref{thm:floor_from_infinitude}, which
itself uses \cref{ax:noncompletable}: were $\mathcal{M}$
completable, the comparison of $u$ against the rest could finish
and the residue could vanish, so realisability with positive cost
implies \cref{ax:noncompletable}.

$(\mathrm{ii}) \Rightarrow (\mathrm{iii})$: This is the body of the
paper. The Floor from Infinitude
(\cref{thm:floor_from_infinitude}) and \cref{thm:floor} give the
floor; \cref{thm:residue_additive,cor:res_accumulate} give the
residue; \cref{thm:category} gives the small-category structure;
\cref{thm:partiality_forced} gives the partial composition;
\cref{thm:residue_cells} gives the residue-cell quotient as the
natural object; \cref{thm:scale_homomorphism} gives the four-type
taxonomy. The Operational Structure Theorem
(\cref{thm:opstructure}) packages these into the quintuple
$\OpStr_\Agent$.

$(\mathrm{iii}) \Rightarrow (\mathrm{i})$: The small category
$\mathbf{C}_{\Agent,\Graph}$ has at least the identity morphisms;
under the floor $\floor > 0$ and the realisability marker, any
hom-set $\mathrm{Hom}(u, v)$ with $u \neq v$ that is non-empty
exhibits a realised non-identity morphism, i.e. a realisable
distinguishability act.

The cycle is closed. $\qed$
\end{proof}

\begin{corollary}[Identifiability is the operational structure]
\label{cor:identifiability_is_structure}
The bare fact that the agent can distinguish any one part of the
system from any other in the ambient whole is sufficient to force
the entire operational structure. Conversely, dissolving the
structure (no floor, total composition, antisymmetric residue
order, or fewer than four operation types) requires denying the
agent any single act of distinguishing parts.
\end{corollary}

\begin{proof}
The forward direction is the cycle of \cref{thm:master}: (i) forces
(ii) forces (iii). The contrapositive is the corollary's second
clause: failure of any clause of (iii) propagates back to failure
of the conjunction in (ii) and thence to absence of any realisable
distinguishing act in (i). $\qed$
\end{proof}



\begin{figure*}[!htbp]
    \centering
    \includegraphics[width=\textwidth]{figures/panel_09_master.png}
    \caption{\textbf{Master Theorem: Identifiability and Operational Structure are Equivalent.}
    (\textbf{A})~Realised acts on the linear $y$-axis ($\sim 170$ to $\sim 215$) versus floor $\beta$ on a log $x$-axis ($10^{-3}$ to $10^{1}$). Four viridis-coloured curves correspond to four fixed graph topologies of edge density $p \in \{0.15, 0.30, 0.45, 0.65\}$, with $\beta$ swept as a uniform edge-weight scale factor. Each curve is constant in $\beta$ across five decades: the realised-act count is a structural property of the graph, independent of the floor. This is the empirical content of Corollary~\ref{cor:identifiability_is_structure}: identifiability is a property of $(\mathcal{A}, \mathcal{G})$, while $\beta$ controls only the cost of realising it.
    (\textbf{B})~Mean realised acts (averaged over $12$ random replicas per $\beta$) on the linear $y$-axis ($0$ to $\sim 90$) versus uniform floor $\beta$ on the linear $x$-axis ($0$ to $1.0$). Steelblue circles trace the mean; the crimson dashed reference at $\beta = 0$ marks the phase boundary. The mean drops abruptly from $\sim 80$ at $\beta > 0$ to exactly $0$ at $\beta = 0$. Empirical content of the contrapositive of Theorem~\ref{thm:master}: setting $\beta = 0$ (denying non-completability, by Corollary~\ref{cor:floor_trace}) dissolves identifiability; no realisable distinguishability act exists.
    (\textbf{C})~Cumulative residue on a log $y$-axis ($10^{-1}$ to $10^{2}$) versus $\beta$ on a log $x$-axis ($10^{-3}$ to $10^{1}$) for the fixed-topology sweep. Steelblue circles trace the measured cumulative residue $\sum\rho_i$; the crimson dashed reference line $\sum\rho_i \propto \beta$ has slope $1$ in log-log coordinates. Measured and theoretical curves coincide across five decades, confirming that cumulative residue scales linearly with the floor and the realised-act count is structural.
    (\textbf{D})~Three-dimensional identifiability phase curve: $\log_{10}\beta$ on the linear $x$-axis ($-3$ to $1$), realised acts on the linear $y$-axis ($\sim 200$ to $\sim 220$), $\log_{10}\sum\rho_i$ on the vertical axis ($-1$ to $4$), viridis colourmap by $\log_{10}\beta$, viewing angle elevation $30^{\circ}$ azimuth $-60^{\circ}$. The curve is monotone in $\log\beta$ and $\log\sum\rho_i$ while flat in realised-act count: the structure is invariant in the agent-graph topology, and the cost surface rises linearly with the floor. Geometric realisation of the three-way equivalence (i) $\Leftrightarrow$ (ii) $\Leftrightarrow$ (iii) of Theorem~\ref{thm:master}.}
    \label{fig:panel-master}
\end{figure*}

% =====================================================================
\section{Implications}
\label{sec:implications}
% =====================================================================

We draw three implications of the main theorem for engineering
practice. Each is a direct consequence of \cref{thm:opstructure} or
its corollaries.

\subsection{Software architecture under the conditional}

The conjunction of \cref{ax:graph} and \cref{ax:agent} covers all
realised software systems: any program with finite memory and
finite execution is a finite agent, and any data structure it
manipulates is a finite graph. By \cref{thm:opstructure}, the
program's operational structure is forced by its data structure and
its agency: it is not chosen.

In practice this means: arguments about which operations a program
``should'' admit reduce to arguments about $(\Agent, \Graph)$.
Two architectures disagreeing about the operational interface are
disagreeing about one of: which finite data structure the system is
modelled as; which finite agent the program is. Once those are
agreed, the operations are determined.

\subsection{Database query algebras}

A finite database is a finite graph in the sense of
\cref{def:graph}: the schema gives the vertices (tables, attributes,
values), the foreign-key relations give the edges, and the
cardinalities give the edge measures. A finite query engine is a
finite agent in the sense of \cref{def:agent}: bounded internal
state, finite query plan space, finite cost budget per query.

By \cref{thm:opstructure}, the query algebra is forced. The choice
between relational algebra, calculus, and CRUD is the choice of
which subset of the forced algebra to expose; it is not the choice
of which algebra to invent. The minimal operational floor is the
smallest positive edge weight in the schema graph (typically a unit
cardinality), and every realisable query has residue $\geq$ this
floor.

\subsection{Knowledge graphs and reasoning}

A finite knowledge graph---ontology, RDF triple store, labelled
property graph---is a finite graph in the sense of \cref{def:graph}.
A finite reasoner---SPARQL engine, description-logic prover, graph
neural network---is a finite agent in the sense of \cref{def:agent}.
By \cref{thm:opstructure}, the reasoner's operational structure is
forced by the graph it queries and the resources it has.

The implication is that the apparent diversity of reasoning systems
collapses, under the conditional, to relabelling. Two reasoners
with the same $(\Agent, \Graph)$ admit the same operations up to
naming; their differences are presentational, not structural.
Cross-system calibration (translation of one query language into
another) is guaranteed to exist whenever $(\Agent, \Graph)$ is
shared.

% =====================================================================
\section{Conclusion}
\label{sec:conclusion}
% =====================================================================

\subsection{What we have established}

We have shown that under two finiteness premises---finite engaging
agent (\cref{ax:agent}), finite engaged graph-structured system
(\cref{ax:graph})---the system necessarily admits an operational
structure $\OpStr_\Agent$ that is forced by the premises and unique
given the pair $(\Agent, \Graph)$ together with the agent's
path-choice policy. The structure is a small category with partial
composition, a residue function with a strictly positive floor on
non-identity morphisms, a residue preorder that descends to a
partial order on the residue quotient, a budget-bounded
realisability marker, and a four-type classification (algebraic,
dynamical, computational, internal) that exhausts the realisable
operations.

The result is a conditional architectural claim. It does not
address what reality is in itself; it addresses what structure
follows when the two conditions hold. Where they hold---which is
the case in essentially all realised engineering systems---the
operational structure is determined.

\subsection{What we have not claimed}

We have not claimed that:
\begin{enumerate}[nosep,leftmargin=*]
\item Reality is fundamentally finite. The premises restrict the
agent and the system the agent engages, not the world.
\item All systems are graphs. The premises identify the conditions
under which the result holds; systems failing the conditions are
not addressed.
\item Operational structures are the only structures of interest.
Other structures (continuous dynamics, categorical limits, infinite
processes) may coexist with operational structure; we have not
spoken to them.
\end{enumerate}

\subsection{The narrowness of the result is the point}

A conditional result is sharp because it is narrow. The framework
established here speaks definitively about a specific class of
situations and is silent about others. This is exactly the right
shape for an architectural claim: scope is a feature, and the scope
of \cref{thm:opstructure} is the scope of finite agency engaging
finite graph-structured systems. Within that scope, the operational
structure is forced. Outside it, the framework declines to address
the case.

\subsection{What this enables}

The main practical consequence is that engineering decisions over
$(\Agent, \Graph)$ become discovery tasks rather than invention
tasks. To design a finite system's operational structure is to
identify the unique $\OpStr_\Agent$ that the agent--graph pair
determines. To compare two designs is to ask whether they realise
the same $(\Agent, \Graph)$. To extend a system is to ask how the
extension changes $(\Agent, \Graph)$ and re-derive the structure.

The framework converts what is currently treated as a matter of
architectural taste into a matter of formal derivation. Whether
this is a useful conversion in practice depends on whether the cost
of formalising $(\Agent, \Graph)$ for a given system is less than
the cost of debating its operational structure informally. We
expect this to be true at scale---for large systems with many
stakeholders and long lifetimes---and false at small scale, where
the informal debate is cheap.

\subsection{Closing}

The structure of operations on a finite graph engaged by a finite
agent is, in its uniquely-determined part, not a design problem. It
is a derivation from the two finiteness premises. The derivation
produces an answer unique up to labelling and path-choice: a small
category with a positive floor, a residue preorder, partial
composition under a budget constraint, and a four-type
classification of effects. That is what the conditional
architecture says, and that is all it says.

% =====================================================================
\begin{thebibliography}{9}

\bibitem[Bollob\'as(1998)]{bollobas1998} B\'ela Bollob\'as.
\emph{Modern Graph Theory}.
Graduate Texts in Mathematics 184, Springer, 1998.

\bibitem[Davey and Priestley(2002)]{daveypriestley2002}
B.~A. Davey and H.~A. Priestley.
\emph{Introduction to Lattices and Order}, 2nd edition.
Cambridge University Press, 2002.

\bibitem[Diestel(2017)]{diestel2017} Reinhard Diestel.
\emph{Graph Theory}, 5th edition.
Graduate Texts in Mathematics 173, Springer, 2017.

\bibitem[Halmos(1950)]{halmos1950} Paul R. Halmos.
\emph{Measure Theory}.
Van Nostrand, 1950.

\bibitem[Hopcroft, Motwani and Ullman(2007)]{hopcroft2007}
John E. Hopcroft, Rajeev Motwani, and Jeffrey D. Ullman.
\emph{Introduction to Automata Theory, Languages, and Computation},
3rd edition. Addison-Wesley, 2007.

\bibitem[Howie(1995)]{howie1995} John M. Howie.
\emph{Fundamentals of Semigroup Theory}.
London Mathematical Society Monographs 12, Oxford University Press,
1995.

\bibitem[Mac Lane(1998)]{maclane1998} Saunders Mac Lane.
\emph{Categories for the Working Mathematician}, 2nd edition.
Graduate Texts in Mathematics 5, Springer, 1998.

\end{thebibliography}

\end{document}
